\documentclass[11pt,twoside]{article}

\usepackage[margin=1in]{geometry}
\usepackage{amsmath,amssymb,amsthm,mathtools}
\usepackage{thmtools}
\usepackage[inline,shortlabels]{enumitem}
\usepackage{booktabs}
\usepackage{xcolor}
\usepackage{microtype}

% Revision date: maintained by scripts/build_lecture_notes.py. Keep this command.
\newcommand{\lecturerevised}{2026-09-15 23:21 UTC-06:00}
\usepackage{eso-pic}
\AddToShipoutPictureBG{%
  \AtPageLowerLeft{%
    \put(\LenToUnit{0.5\paperwidth},\LenToUnit{0.5in}){%
      \makebox(0,0){\normalfont\footnotesize\color{black!60}Last revised: \lecturerevised}%
    }%
  }%
}
\usepackage[square,authoryear]{natbib}
\usepackage{xr-hyper}
\usepackage[colorlinks=true,linkcolor=blue!60!black,
  urlcolor=blue!60!black,citecolor=blue!60!black]{hyperref}
\externaldocument[lecture2-]{lecture02}
\externaldocument[lecture3-]{lecture03}
\usepackage[nameinlink,capitalize]{cleveref}
\hypersetup{pdftitle={Lecture 4: Sequence modeling, generation, and memory},
  pdfauthor={Csaba Szepesvari}}
\crefname{exercise}{Exercise}{Exercises}
\Crefname{exercise}{Exercise}{Exercises}

% Short refinements that may be skipped on a first reading.
\usepackage{enotez}
\DeclareInstance{enotez-list}{lecture}{list}{format=\normalsize,number={#1.}}
\setenotez{list-name=Endnotes,list-style=lecture,backref=true}
\crefname{endnote}{Endnote}{Endnotes}
\Crefname{endnote}{Endnote}{Endnotes}

\definecolor{courseblue}{HTML}{17365D}
\newcommand{\R}{\mathbb R}
\newcommand{\E}{\mathbb E}
\newcommand{\I}{\mathbb I}
\renewcommand{\P}{\mathbb P}
\newcommand{\KL}{D}
\newcommand{\TV}{\mathrm{TV}}
\newcommand{\kl}{\mathrm{kl}}
\newcommand{\tv}[1]{\lVert #1\rVert_{\mathrm{TV}}}
\newcommand{\suppmark}{\texorpdfstring{\textsuperscript{*}}{*}}
\DeclareMathOperator{\diag}{diag}
\DeclareMathOperator{\rank}{rank}
\makeatletter
\@for\@letter:=A,B,C,D,E,F,G,H,I,J,K,L,M,N,O,P,Q,R,S,T,U,V,W,X,Y,Z\do{%
  \expandafter\edef\csname c\@letter\endcsname{\noexpand\mathcal{\@letter}}}
\makeatother
\newtheorem{theorem}{Theorem}[section]
\newtheorem{proposition}[theorem]{Proposition}
\newtheorem{corollary}[theorem]{Corollary}
\theoremstyle{definition}
\newtheorem{definition}[theorem]{Definition}
\newtheorem{example}[theorem]{Example}
\theoremstyle{remark}
\newtheorem{remark}[theorem]{Remark}
\renewcommand{\thepage}{4--\arabic{page}}
\renewcommand{\thesection}{4.\arabic{section}}
\renewcommand{\theequation}{4.\arabic{equation}}
\pagestyle{myheadings}
\markboth{Lecture 4: Sequence modeling, generation, and memory}
{Lecture 4: Sequence modeling, generation, and memory}
\usepackage{../exercises/course-exercises}

\begin{document}

\begin{center}
\fcolorbox{courseblue}{white}{%
\begin{minipage}{0.92\textwidth}
\vspace{1mm}
\textbf{CMPUT 654: Mathematical Foundations of Modern AI Systems}
\hfill Fall 2026
\vspace{3mm}
\begin{center}
{\Large\bfseries Lecture 4: Sequence modeling, generation, and memory}
\vspace{1mm}

September 10, 2026
\end{center}
\vspace{2mm}
\textit{Lecturer: Csaba Szepesv\'ari}
\hfill
\textit{Note prepared by: Csaba Szepesv\'ari}
\vspace{1mm}
\end{minipage}}
\end{center}

\noindent\textbf{About these lecture notes.}
We know how to fit a sequence model by next-token prediction. We now ask
whether fitting successful examples makes the model generate successful
sequences. The answer depends on how prediction errors affect later
generation, what cost we care about, and what structure the model preserves.
This leads from log loss to error recovery, reverse KL, and models with
explicit memory.

\medskip
\noindent\textbf{Scope relative to class.}
Further questions and ideas arose while preparing these notes, so they
extend and refine the discussion from the lecture.
The lecture developed the positive-data fitting question, the binary
error-propagation example, stability of Markov models, and the
concentration-dependent reverse-KL bound. The state-space introduction
and system-identification material are now collected in Lecture~5,
which develops these topics in detail.
The notes group the cost bounds first, then use exact fixed-KL cost
intervals to develop two laws attaining the cost bounds. We reuse these
laws when interpreting normalized loss, comparing fitting objectives,
and introducing recovery. Repeated errors, parity, and the explicit
modulus calculations are developed in exercises.
A further section asks how parameter sharing affects learning as
sequence length grows. The fitting and structural choices follow.
Headings ending in a superscript asterisk\suppmark{} mark supplementary
material: sharper comparisons
between likelihood and cost and a finite-sample fitting guarantee. The exercises retain the calculations and proofs left to
students. Longer proofs, a statistical obstruction to reverse-KL fitting,
reverse-KL optimization details, and recovery for history-dependent generators
are in the appendices.

\medskip
\noindent\textbf{Notation and reading guide.}
Tokens are \(X_t\in[N]\). We use \(S_t\) for a state when
discussing persistent and revisable memory.
A vector $x\in \R^d$ is a column vector.
For indices \(a\le b\), write \(X_{a:b}=(X_a,\ldots,X_b)\) and
\(X_{<t}=X_{1:t-1}\); lowercase \(x_{a:b}\) denotes a realized sequence.
For a symbol \(u\), \(u^k=(u,\ldots,u)\) is the length-\(k\) constant
sequence, so \(0^T=(0,\ldots,0)\), and juxtaposition denotes
concatenation, as in \(10^{T-1}=(1,0,\ldots,0)\).
Sequence notation is useful because the same object has two roles: it is an
ordered tuple in \([N]^T\), and it is a time-indexed function on \([T]\).
We use tuple notation when order and prefixes matter, and ordinary vector
notation when applying linear algebra.
For a sequence cost \(c: [N]^T\to \R\) and a distribution \(P\) over
\([N]^T\), \(\E_P[c]\) abbreviates \(\E[c(X_{1:T})]\) for
\(X_{1:T}\sim P\). Thus
\[
  \E_P[c]=\sum_{x_{1:T}}P(x_{1:T})c(x_{1:T}).
\]
By default, in what follows, \(P\) is the data law and \(Q\), or \(Q_\theta\), is
a model law.
We use \(\kappa\) for a KL value or budget, \(s\) for a cost-discrepancy
value or budget, and \(\eta\) for per-token accuracy, including the
optimization tolerance introduced in \cref{sec:likelihood-pitfall}.
Logarithms are natural. Read through \cref{sec:take-home}; the supplementary
material can be read separately. Numbered endnotes, bibliographical notes and
references, the glossary, exercises, and appendices follow the take-home
messages in that order.

\section{From next-token prediction to useful generation}
\label{sec:prediction-generation}

Lecture~1 described transformers as maps from token prefixes to
next-token distributions. Lecture~2 asked what structure allows learning
to generalize. Lecture~3 explained why log loss is a natural objective for
probability estimation. Pretraining brings these ideas together: shared
parameters are fitted to predict tokens across many sequences and contexts.
The prediction problem supplies a scalable way to learn from data.

Throughout the discrete discussion, sequences have fixed length \(T\)
over a finite vocabulary \([N]\), and costs are real-valued. The finite
space \([N]^T\) simplifies the arguments: every cost is bounded, and
expectations are finite sums. Many results hold more generally;
\cref{app:finiteness} explains where finiteness is used and what can
replace it. It is also easy to extend the results to allow variable length
sequences as outlines in Lecture 1. For simplicity, we restrict our attention
to the fixed length case.

Let \(X^{(1)}_{1:T},\ldots,X^{(n)}_{1:T}\) be independent training sequences
from \(P\), a distribution over $[N]^T$.
Their tokens may have arbitrary dependence within each sequence.
Empirical log-loss minimization chooses
\begin{equation}
  \widehat\theta\in\arg\min_\theta
  \frac1n\sum_{i=1}^n-\log Q_\theta(X^{(i)}_{1:T}).
  \label{eq:empirical-likelihood}
\end{equation}

The chain rule of probabilities gives that for every $x_{1:T}\in [N]^T$,
\begin{equation}
  Q_\theta(x_{1:T})
  =\prod_{t=1}^T Q_\theta(x_t\mid x_{<t}),
  \qquad x_{<t}=(x_1,\ldots,x_{t-1}).
  \label{eq:probability-chain}
\end{equation}
Every sequence law has this factorization. We call it the
\emph{autoregressive factorization}; restrictions on the conditional
predictors, such as a transformer architecture or a finite context window,
specify the model class. Taking logarithms in
\cref{eq:probability-chain} and plugging in $X_{1:T}$ in place of
$x_{1:T}$ gives
\begin{equation}
  -\log Q_\theta(X_{1:T})
  =\sum_{t=1}^T-\log Q_\theta(X_t\mid X_{<t}).
  \label{eq:next-token-loss}
\end{equation}
Consequently, next-token log-loss minimization, the workhorse of language model training,
 is sequence maximum likelihood.

When evaluating \cref{eq:next-token-loss}, the predictor receives the
observed prefix, including the observed next input.
This is \emph{teacher forcing}:
the student's cost is measured at the
teacher's data. The metaphore comes from viewing the problem as a teacher generating
data that a student needs to use to learn.
During generation, the
next input is instead sampled from the model:
\begin{equation}
  \widetilde X_t\sim
  Q_\theta(\,\cdot\mid\widetilde X_{<t}).
  \label{eq:generation}
\end{equation}
The resulting sequence has joint law \(Q_\theta\). The same factorization
supports both likelihood evaluation and sampling.


As explained earlier,
the parameter $\widehat \theta$ from
\cref{eq:empirical-likelihood} is
maximum likelihood parameter estimate:
minimizing the negative log of a product
maximizes that product, which, under independence, is the probability assigned to the data.
For a fresh sequence \(X_{1:T}\) with law \(P\),
the log-loss identity from Lecture~3 and the decomposition in
\cref{eq:next-token-loss} give
\begin{equation}
  \begin{aligned}
  \E[-\log Q_\theta(X_{1:T})]
  &=H(P)+\KL(P, Q_\theta),\\
  \KL(P, Q_\theta)
  &=\sum_{t=1}^T\E\!\left[
    \KL\bigl(P(\,\cdot\mid X_{<t})
       , Q_\theta(\,\cdot\mid X_{<t})\bigr)\right].
  \end{aligned}
  \label{eq:population-kl}
\end{equation}
The entropy term is independent of \(\theta\). Thus minimizing population
log loss minimizes KL; empirical log loss estimates this population
objective from data.
The second identity is the \emph{KL chain rule}: sequence KL is the
sum of the expected conditional KLs. \Cref{ex:kl-chain} asks for its
derivation.





In ML/statistics we distinguish between data and model laws.
This leads to the following definition:
\begin{definition}[Forward and reverse KL]
When  \(P\) is a data law and \(Q\) is a model law, we call
$\KL(P, Q)$ the \emph{forward KL}, while
$\KL(Q, P)$ is called the \emph{reverse KL}.
\end{definition}
With our finite spaces, we have
\[
  \begin{aligned}
  \KL(P, Q)&=\sum_xP(x)\log\frac{P(x)}{Q(x)}
    &&\text{(the \emph{forward KL})},\\
  \KL(Q, P)&=\sum_xQ(x)\log\frac{Q(x)}{P(x)}
    &&\text{(the \emph{reverse KL})}.
  \end{aligned}
\]
Note that above a zero numerator contributes zero;
a positive numerator with zero denominator contributes \(+\infty\).\endnote{\textbf{Conventions at zero.}
These conventions are the continuous extensions of the nonzero formula and
express the infinite log loss incurred when a positive-probability outcome is
assigned probability zero; Lecture~3, \cref{lecture3-ex:proper-regrets}(e),
verifies both limits.}
The names refer to the roles of \(P,Q\): forward KL averages under the
data law, whereas reverse KL averages under the generated law.

As explained, $\widehat\theta$, is expected to approximately minimize $\KL(P, Q_\theta)$ in $\theta$.
The discrepancy $\KL(P, Q_{\widehat\theta})- \inf_{\theta}\KL(P, Q_\theta)$
describes the effect of using a finite dataset in the lack of access to $P$ (and a reliable way of minimizing $\KL(P, Q_\theta)$).
Recall from Lecture 2 that this is expected to be controlled when $n$ is large and things are ``sufficiently nice''.
As such, we leave this question alone in this lecture and instead ask, would it be even
beneficial (and to what degree) if we could find a model $Q$ with a small population KL?


\subsection{A task specifies which generations are useful}

A prompt asks the model to continue in a particular way. Mathematically,
\emph{prompting} uses a conditional distribution
\(Q_\theta(\,\cdot\mid u)\), where \(u\) is the supplied prompt
(we omit the full formal description of the setup, but here $u$
is also a sequence in the typical language model training setting).
For a mathematics task, a useful continuation might be a correct solution;
for a knowledge question, an accurate answer. A \emph{cost}
\(c(u,x_{1:T})\) assigns a numerical value to the generated response.
Given a test-prompt distribution \(\nu\),
the expected cost of a model $Q_\theta$ is
\begin{equation}
  \mathcal J(\theta)=
  \E[c(U,\widetilde X_{1:T})],
  \qquad U\sim\nu,\quad
  \widetilde X_{1:T}\mid U\sim Q_\theta(\,\cdot\mid U).
  \label{eq:prompt-cost}
\end{equation}
A model's quality then is determined by $\mathcal J(\theta)$.
The cost function is thus $c: [N]^T \to \R$.
In what follows, without loss of generality,
we assume that
it is nonnegative valued.
Sometimes (especially in RL) the task is formulated with rewards.
These formulations are equivalent:
Maximizing expected reward is equivalent to minimizing its negative.

A common approach to get a small-cost model
is to collect low-cost (or even zero-cost) prompt--response pairs,
and then find a model that minimizes the empirical log-loss, aiming to
``mimick'' the responses in the training data.
Formally, solve
\begin{equation}
  \widehat\theta\in\arg\min_\theta
  \frac1n\sum_{i=1}^n\sum_{t=1}^T
  -\log Q_\theta(X^{(i)}_t\mid U^{(i)},X^{(i)}_{<t}).
  \label{eq:finetuning}
\end{equation}
In training large language models, often incremental methods are used, such
as some form of gradient descent.
Lectures~9--11 return to how gradient-based optimization selects among fitted
solutions and when local gradients lead to global convergence.
When this is the case, typically people start gradient descent with
a parameter $\theta_0$ which is obtained during \emph{pre-training} (as outlined in Lecture 1).
Then, with the data as above, people continue gradient descent
from $\theta_0$ with the data as above.
This practice is called  \emph{fine-tuning} by likelihood on \emph{positive examples}
(starting from a pretrained model parameters).
The data is called ``positive'' because it will contain completed proofs (in math problems)
or, more generally, answered deemed correct,
with no examples in the data of incorrect responses.

For simplicity, in this lecture we do not consider the ``continuing'' aspect
of the fine-tuning and will ask the equation whether minimizing log-loss
with positive data will produce low-cost answers to new prompts.
To isolate the connection between fitting and task performance, we set the
prompts aside; keeping the prompts (conditioning) around would introduce
extra notation, but otherwise, in general, changes nothing. There is one
exception to this which will be noted when it arises.

With this, let is turn to formulating the specific question studied next.
Let \(P\) be the law of positive sequences, \(Q\) the
fitted law, and \(c\) the cost. One should generally think that \(\E_P[c]=0\);
the results will aim at this case but will be more general.
Recall from \cref{eq:population-kl} that the population log loss is
\[
 \E_P[-\log Q]=H(P)+\KL(P, Q).
\]
As explained earlier, one should think about log-loss minimization, or
likelihood fitting as seeking a model $Q$ that with small $\KL(P, Q)$.
Yet, the task calls
for small generated cost \(\E_Q[c]\).
Suppose fitting gives us \(\KL(P, Q)\le\kappa\).
The basic question we ask is how large can \(\E_Q[c]\) be?
In other words, is controlling KL this way sufficient to control generation
cost? To what degree? What are the strengths and weaknesses of this approach?
Naturally, we will continue with asking whether there is anything better.

\section{How does KL control the task cost?}
\label{sec:kl-cost}

Define the \emph{cost discrepancy}
\[
 d_c(P,Q)=|\E_P[c]-\E_Q[c]|
\]
to measure the change in expected cost from the data law to the model law.
For cost-free data laws,
thank to the assumption that $c\ge 0$,
\(d_c(P,Q)=\E_Q[c]\), so a bound on this
discrepancy answers our question.

We will compare four ways to turn KL control into cost control, starting with simpler
and moving to more elaborate results.
We will compare the resulting bounds, and ask which of them are sharp and when.
A supplementary subsection asks for both the smallest and largest cost
discrepancies at a fixed KL, relating these to moduli of continuity.
We then compute these intervals in binary examples to see how differently
models with the same likelihood can behave.

\subsection{Use the cost range}

Recall the definition of the total variation distance between probability laws:
In what follows we fix a finite set $\cX$ as the domain of the cost and the probability distributions $P$ and $Q$.
\begin{definition}[Total variation]
For laws \(P,Q\) over a finite set $\cX$, their
\emph{total-variation distance} is
\begin{equation}
  \tv{P-Q}
  =\sup_{A\subseteq\mathcal X}|P(A)-Q(A)|
  =\frac12\sum_{x\in\mathcal X}|P(x)-Q(x)|.
  \label{eq:tv}
\end{equation}
\end{definition}
Lecture~3 explains why the same formula is a norm on signed measures and a
metric on probability laws, and contrasts both with a pseudometric.
\begin{proposition}[Total variation controls bounded costs]
\label{prop:tv-cost}
If \(0\le c(x)\le C_{\max}\) for every \(x\in \cX\), then
\begin{equation}
  |\E_Q[c]-\E_P[c]|
  \le C_{\max}\tv{P-Q}.
  \label{eq:tv-cost}
\end{equation}
\end{proposition}

This is Lecture~3, \cref{lecture3-ex:pinsker}(c), after scaling the cost by
\(C_{\max}\). The exercise proves it from the positive and negative parts of
the signed measure \(Q-P\).

This bound is exact when maximizing over costs with the specified range:
\(c(x)=C_{\max}\I\{Q(x)>P(x)\}\) attains equality.
Thus TV measures the largest expected-cost discrepancy among such costs.

Pinsker's inequality from Lecture~3 now gives
\begin{equation}
  \E_Q[c]\le
  \E_P[c]+C_{\max}\sqrt{\frac{\KL(P, Q)}2}.
  \label{eq:pinsker-cost}
\end{equation}
For any bounded cost, sufficiently small KL therefore
gives sufficiently small cost discrepancy.
If \(\E_P[c]=0\), a sufficient condition for
\(\E_Q[c]\le s\) is
\[
  \KL(P, Q)\le 2(s/C_{\max})^2.
\]
Both the square-root dependence and its constant are necessary for
general bounded costs.
\begin{example}[A locally sharp fair-coin comparison]
\label{exmp:fair-coin}
Let $\cX=\{0,1\}$,
\[
 P=\operatorname{Bernoulli}(1/2),\qquad
 Q=\operatorname{Bernoulli}(1/2+\delta),\qquad c(x)=x.
\]
As \(\delta\to0\),
\[
 d_c(P,Q)=|\delta|,\qquad
 \KL(P, Q)=2\delta^2+O(\delta^4).
\]
Pinsker's bound divided by the actual discrepancy tends to one: it is
\emph{locally sharp}.
\end{example}
For a count of token errors, \(C_{\max}=T\),
so the horizon also enters the conversion.

Recall that we are interested in the case when $\E_P c=0$ or small, while $c\ge 0$.
Can we improve this bound using this information?

\subsection{Successful data give a sharper bound\suppmark}

Let us start with the case when $\E_P c=0$.




Then, the following result applies:

\begin{proposition}[Forward KL and zero-cost data]
\label{prop:positive-data}
If \(0\le c\le C_{\max}\), \(\E_P[c]=0\)
then
\begin{equation}
 \E_Q[c]\le C_{\max}(1-e^{-\KL(P, Q)})
 \le C_{\max}\KL(P, Q).
 \label{eq:positive-data-bound}
\end{equation}
\end{proposition}
This is an improvement over \cref{eq:pinsker-cost}
when $\KL(P, Q)\le 1/2$ and the gap between them grows as the KL between $P$ and $Q$ approaches zero.
For small KL, exploiting successful data replaces the
square-root dependence from Pinsker by linear dependence.


Recording only whether an event occurred can only decrease KL. For an event certain under \(P\), this gives a lower bound on its probability under \(Q\).
Recall that for $p,q\in [0,1]$, $\kl(p,q)$ denotes the KL between the Bernoulli distributions with parameters $p$ and $q$.
\begin{proposition}[From distribution KL to event probability KL]
Let $P,Q$ be distributions over the finite domain $\cX$.
Then, for any $A\subset \cX$,
\[
\kl(P(A),Q(A))\le \KL(P, Q)\,.
\]
In particular, if $P(A)=1$ then
\[
Q(A) \ge e^{-\KL(P, Q)}\,.
\]
\end{proposition}
\begin{proof}
Consider the transformation that maps
any $x\in \cX$ to $1$ if $x\in A$ and to $0$ otherwise (this is just the indicator if $A$).
Apply the data processing inequality with this deterministic transformation to get the first inequality.
Indeed, the push-forward of $P$ ($Q$) under $\I_A$ is the Bernoulli distribution
with parameter $P(A)$ (respectively, $Q(A)$).
To get the second inequality, note that if $P(A)=1$, $\kl(P(A),Q(A)) = \log \frac{P(A)}{Q(A)}$.
Rearrangement then gives the second inequality. The general data-processing
statement and its log-sum proof appear in Lecture~3,
\cref{lecture3-prop:kl-data-processing,lecture3-ex:pinsker}.
\end{proof}
With this we can return to the proof of \cref{prop:positive-data}.
\begin{proof}[Proof of \cref{prop:positive-data}]
Let \(A=\{x:c(x)=0\}\). Note that since $c\ge 0$ and $\E_P[c]=0$, $P(c=0)=P(A)=1$.
Now, by algebra, using the result of the previous proposition and since $1-e^{-u}\le u$ for $u\ge 0$,
\[
\E_Q \, c
\le
C_{\max} Q(A^c)
=
C_{\max} (1-Q(A)) \le C_{\max}(1-e^{-\KL(P, Q)})
\le C_{\max} \KL(P, Q)
\,. \qedhere
\]
\end{proof}

The bound is sharp:
Let $c$ take on values of $0$ and $C_{\max}$, but otherwise be arbitrary.
Let $P$ be as in the proposition.
Let $R$ be any probability distribution supported on
$\{c>0\} = \{ c = C_{\max} \}$.
For any \(0<a<1\), choosing $Q = (1-a) P + a R$, we have
$\KL(P, Q)=-\log(1-a)$
and
$\E_Q[c]=aC_{\max}$.
Hence, indeed
\[
 \E_Q[c]=aC_{\max} = C_{\max} (1- e^{-\KL(P, Q)})\,.
\]
Thus, equality holds in the first bound of
\cref{eq:positive-data-bound}.

\subsection{Concentration and the direction of KL}
\label{sec:cost-concentration}
Our next result strengthens the bound we obtained
using Pinsker inequality when
the costs under $P$ concentrate around their mean,
but the price is switching from forward to reverse KL:

\begin{restatable}[Cost control from exponential moments]{proposition}{entropycost}
\label{prop:entropy-cost}
Let $P,Q$ be distributions over a finite set $\cX$ and let $c:\cX \to \R$. Then.
\begin{equation}
 \E_Q[c]-\E_P[c]
 \le\inf_{\lambda>0}
 \frac{\KL(Q, P)
       +\log\E_P[e^{\lambda(c-\E_P[c])}]}{\lambda}.
 \label{eq:entropy-cost}
\end{equation}
\end{restatable}

This reduces cost control to bounding an exponential moment and
optimizing one positive scalar. The exponential-tilting proof is in
\cref{app:reverse-proof}. A quadratic bound on that logarithmic
moment function gives a particularly simple answer.

\begin{definition}[Subgaussian cost]
The cost \(c\) is \emph{\(v^2\)-subgaussian under \(P\)}, for \(v\ge0\),
if
\begin{equation}
 \log\E_P\!\left[
 e^{\lambda(c-\E_P[c])}\right]
 \le\frac{\lambda^2v^2}{2}
 \quad\text{for every }\lambda\in\R.
 \label{eq:subgaussian}
\end{equation}
\end{definition}

Here \(v\) measures a scale of fluctuations; \(v^2\) is the variance
proxy.
We write \(v_P(c)\) for the \emph{optimal variance proxy},
which is also known
as the subgaussian constant of $c$ under $P$:
\[
 v_P(c)=\min\left\{a\ge0:
   \log\E_P[e^{\lambda(c-\E_P[c])}]
   \le\frac{\lambda^2a}{2}\ \text{for every }\lambda\in\R\right\}.
\]
The minimum is attained; a constant cost has \(v_P(c)=0\).\endnote{\textbf{Variance proxy and variance.}
How does the optimal proxy compare with the variance? Write
\(\sigma_P(c)^2=\E_P[(c-\E_P[c])^2]\). Then
\[
 \sigma_P(c)^2\le v_P(c)\le\frac{C_{\max}^2}{4}
 \qquad\text{when }0\le c\le C_{\max}.
\]
The lower bound follows by expanding \cref{eq:subgaussian} at
\(\lambda=0\); the upper bound is Hoeffding's lemma.
Thus \(\sqrt{v_P(c)}\), the subgaussian fluctuation scale, lies between
\(\sigma_P(c)\) and \(C_{\max}/2\).
For rare errors the proxy can greatly exceed the variance;
\cref{app:bound-comparison} gives a Bernoulli example.
The bounds and existence of the minimum are proved in
\cref{app:optimal-proxy}.}


\begin{theorem}[Reverse KL and cost concentration]
\label{thm:reverse-cost}
Suppose \(c\) satisfies \cref{eq:subgaussian}.
Then, interpreting $0\cdot \infty = \infty$, we have
\begin{equation}
 |\E_Q[c]-\E_P[c]|
 \le v\sqrt{2\KL(Q, P)}.
 \label{eq:reverse-cost}
\end{equation}
\end{theorem}
\begin{proof}
Substituting \cref{eq:subgaussian} into \cref{eq:entropy-cost} gives
\[
 \inf_{\lambda>0}
 \left(\frac{\KL(Q, P)}{\lambda}+\frac{\lambda v^2}{2}\right)
 =v\sqrt{2\KL(Q, P)}.
\]
Applying the same argument to \(-c\) bounds the other direction.
\end{proof}

Interchanging \(P,Q\) and using the optimal proxies \(v_P(c),v_Q(c)\)
gives
\begin{equation}
 d_c(P,Q)\le
 \sqrt{2\min\{v_Q(c)\KL(P, Q),\,
                 v_P(c)\KL(Q, P)\}}.
 \label{eq:two-concentration-directions}
\end{equation}
While the forward KL term refines the bound coming from
Pinsker's inequality (i.e., \cref{eq:pinsker-cost}),
when besides that $P$ is supported on zero-cost sequences
no cost information is available while fitting $Q$,
the usefulness of this refinement is limited: $v_Q(c)$ may be as large
as $C_{\max}^2/4$ and remains uncontrolled.

The reverse KL part of the bound
is more promising in that $v_P(c)$ is
something we expect to be small.
In particular, when $P(c=0)=1$, as noted beforehand,
$v_P(c)=0$. However, this immediately reveals a limitation:
The bound tells us that for such $P$ and $c$,
the costs under $P$ and $Q$ either
match perfectly (when $\KL(Q, P)<\infty$)
or if the costs are different then $\KL(Q, P)=\infty$.
This foreshadows the difficulty of controlling reverse KL,
which we will discuss later.

\subsection{Use the variability of the data's cost\suppmark}
\label{sec:variance-comparison}

Can we use cost variation under \(P\) while keeping forward KL,
the quantity likelihood fitting controls?
As it turns out, the answer is yes. The proof of this is slightly more complicated
than the previous ones.
The proof centers the cost under \(P\), expands \(Q-P\) in powers of
\(\sqrt Q-\sqrt P\), and uses Cauchy--Schwarz to obtain a variance
term instead of an exponential-moment bound.
Write \(\sigma_P(c)^2=\E_P[(c-\E_P[c])^2]\) for the variance of the cost.

\begin{restatable}[Cost variance and forward KL]{proposition}{forwardvariance}
\label{prop:forward-variance}
If \(0\le c\le C_{\max}\) and \(\KL(P, Q)<\infty\), then
\begin{equation}
 d_c(P,Q)\le2\sigma_P(c)\sqrt{\KL(P, Q)}
             +C_{\max}\KL(P, Q).
 \label{eq:forward-variance}
\end{equation}
\end{restatable}
\Cref{app:hellinger} gives the proof.

The combination of a variance-dependent square-root term and a
range-dependent linear term is called \emph{Bernstein form}.
The square-root term uses the data's cost variation rather than its
full range.
When \(\sigma_P(c)=0\), this term vanishes, leaving a bound
linear in KL.
Note that this bound extends
the zero-cost data bound \cref{eq:positive-data-bound}.
Indeed, when $c=0$ $P$-almost surely,
$\sigma_P(c)=0$ and we recover the linear bound in
\cref{eq:positive-data-bound}.
The additive term accounts for changes in cost on outcomes
to which \(Q\) assigns probability, including outcomes absent under \(P\).

The variance-sensitive and subgaussian bounds are useful in different
regimes. In \cref{exmp:fair-coin}, this bound is asymptotic to
\(\sqrt2|\delta|\), while Pinsker and the subgaussian bound are
asymptotic to the actual discrepancy \(|\delta|\).
For rare errors, the variance-sensitive bound can instead give the
correct order while the subgaussian bound is much larger;
\cref{app:bound-comparison} works out both comparisons.

\subsection{The play of cost when forward KL is controlled}
\label{sec:discrepancy-comparison}

A sharp upper bound tells us how large the cost discrepancy can be.
To understand how much a KL value tells us about task performance, we
need both the smallest and the largest cost discrepancies compatible
with that value. The binary examples below will give exact intervals:
at the same KL, some generators have much smaller cost than others.
Calculating these fixed-KL intervals reveals the ambiguity that remains
even after the upper bounds have been made sharp.
If the interval is large, we can conclude that KL on its own
does allow a ``large play''.
\emph{Knowing how large the play is exactly what allows us to think
about whether it is worthwhile to go beyond controlling KL.}
The rest of this section develops these ideas and connects them
to standard concepts in mathematics.

Fix the data law \(P\) and cost \(c\), and let \(Q\) range over all
probability laws on \(\cX\). Define the two endpoints
\begin{equation}
 \begin{split}
 f_P(\kappa)&=\min_{Q:\,\KL(P, Q)=\kappa}d_c(P,Q),\\
 g_P(\kappa)&=\max_{Q:\,\KL(P, Q)=\kappa}d_c(P,Q).
 \end{split}
 \label{eq:fixed-kl-endpoints}
\end{equation}
Here KL levels range over the attainable values in \([0,+\infty]\).
The extrema exist in our finite setting; \cref{app:finiteness}
gives the compactness argument. These are the sharp bounds in
\[
 f_P\bigl(\KL(P, Q)\bigr)
 \le d_c(P,Q)\le
 g_P\bigl(\KL(P, Q)\bigr).
\]
The preceding inequalities studied the upper endpoint. The lower
endpoint answers an equally relevant question: \emph{how small could
the cost discrepancy be at this same KL?} A large gap between the
endpoints means that likelihood alone leaves substantial uncertainty
about task performance.
Conversely, a small gap means no such wiggle room,
which would mean that controlling KL alone could be sufficient
to get a good control on the task cost.

\paragraph{From fixed values to error budgets.}
Usually we know an upper bound on a discrepancy, rather than its exact
value. The corresponding sharp conversion functions are called
\emph{moduli of continuity}. At the fixed reference \(P\), define
\begin{equation}
 \omega_{D\to c}(\kappa;P)
 =\sup_{Q:\,\KL(P, Q)\le\kappa}d_c(P,Q).
 \label{eq:cost-modulus}
\end{equation}
This is the modulus of continuity of the cost discrepancy with respect to KL. It asks:
\emph{how large can the cost discrepancy be when KL is at most \(\kappa\)?}
In the other direction, define
\begin{equation}
 \omega_{c\to D}(s;P)
 =\sup_{Q:\,d_c(P,Q)\le s}\KL(P, Q).
 \label{eq:kl-modulus}
\end{equation}
This is the modulus of continuity of KL with respect to cost discrepancy. It asks:
\emph{how large can KL be when the cost discrepancy is at most \(s\)?}
The arrows specify which discrepancy controls the other; both formulas
use forward KL, \(\KL(P, Q)\).

\begin{proposition}[The sharp conversions from fixed-KL endpoints]
\label{prop:endpoints-moduli}
For the fixed \(P,c\),
\begin{equation}
 \begin{split}
 \omega_{D\to c}(\kappa;P)&=\sup_{\kappa'\le\kappa}g_P(\kappa'),\\
 \omega_{c\to D}(s;P)&=\sup\{\kappa:f_P(\kappa)\le s\}.
 \end{split}
 \label{eq:endpoints-moduli}
\end{equation}
Both moduli are nondecreasing. They are the smallest nondecreasing
functions giving, respectively, the bounds
\[
 d_c(P,Q)\le\omega_{D\to c}\bigl(\KL(P, Q);P\bigr),
 \qquad
 \KL(P, Q)\le\omega_{c\to D}\bigl(d_c(P,Q);P\bigr).
\]
\end{proposition}

\begin{proof}
At KL level \(\kappa\), the largest cost discrepancy is \(g_P(\kappa)\),
and some cost discrepancy is at most \(s\) exactly when
\(f_P(\kappa)\le s\).
Taking the indicated suprema proves \cref{eq:endpoints-moduli}.
Increasing either budget enlarges its feasible set, proving
monotonicity. If a nondecreasing function \(g\) satisfies
\(d_c(P,Q)\le g(\KL(P, Q))\), then every \(Q\) with KL at most
\(\kappa\) has discrepancy at most \(g(\kappa)\). Taking the supremum proves
\(\omega_{D\to c}(\kappa;P)\le g(\kappa)\). The other direction is identical.
\end{proof}

Thus the upper endpoints determine the \(D\to c\) modulus, while the
lower endpoints determine the \(c\to D\) modulus. In particular, if
\(g_P\) is nondecreasing, it already equals \(\omega_{D\to c}\).
An increasing lower bound can be read in the opposite direction by
inverting it.

\begin{proposition}[Inverting a lower bound]
\label{prop:lower-bound-inversion}
Suppose \(f:[0,+\infty]\to[0,+\infty]\) is strictly increasing and
\(f(\KL(P, Q))\le d_c(P,Q)\) for every \(Q\). Then
\[
 \omega_{c\to D}(f(\kappa);P)\le\kappa.
\]
If some \(Q\) has \(\KL(P, Q)=\kappa\) and \(d_c(P,Q)=f(\kappa)\),
equality holds. In particular, when the sharp lower endpoint \(f_P\)
is strictly increasing,
\(\omega_{c\to D}(s;P)=f_P^{-1}(s)\) for \(s\) in its range.
\end{proposition}

\begin{proof}
If \(d_c(P,Q)\le f(\kappa)\), then
\(f(\KL(P, Q))\le f(\kappa)\). Strict increase gives
\(\KL(P, Q)\le\kappa\). A law attaining the lower bound at \(\kappa\)
also attains this KL budget, proving equality.
\end{proof}

This explains why we care about the converse modulus: it describes
the lower-bound side of the same comparison. A small lower endpoint
and a large upper endpoint allow models of very different task quality
to have exactly the same KL. The examples below will compute both
endpoints and exhibit such models.

\paragraph{Uniform comparisons over data laws.}
If the reference law is unknown, we may want a conversion valid for
every \(P\) in a class \(\mathcal P\). The \emph{uniform moduli} are
\[
 \omega^{\mathcal P}_{D\to c}(\kappa)
   =\sup_{P\in\mathcal P}\omega_{D\to c}(\kappa;P),
 \qquad
 \omega^{\mathcal P}_{c\to D}(s)
   =\sup_{P\in\mathcal P}\omega_{c\to D}(s;P).
\]
Enlarging \(\mathcal P\) can only increase these worst-case conversion
bounds. For example, with \(0\le c\le C_{\max}\), take
\(\mathcal P_\sigma=\{P:\sigma_P(c)\le\sigma\}\).
\Cref{prop:forward-variance} gives
\[
 \omega^{\mathcal P_\sigma}_{D\to c}(\kappa)
 \le 2\sigma\sqrt\kappa+C_{\max}\kappa.
\]
This reformulates the earlier sharpness question: among all data laws
whose cost variance is at most \(\sigma^2\), how much can cost change
under a KL budget \(\kappa\)? Throughout, \(Q\) still ranges over all laws.

\paragraph{The norm-comparison analogy.}
For two norms on \(\R^d\), sharp bounds
\(\alpha\|z\|_a\le\|z\|_b\le\beta\|z\|_a\) place
\(\|z\|_b\) between \(\alpha r\) and \(\beta r\) when
\(\|z\|_a=r\). The ratio \(\beta/\alpha\) measures the
multiplicative spread at this fixed input norm.
For norms the sharp budget conversion is linear in the budget \(r\):
substituting \(z=ru\) gives
\[
 \sup_{\|z\|_a\le r}\|z\|_b
 =r\sup_{\|u\|_a\le1}\|u\|_b
 \qquad(r>0).
\]
Norms scale linearly under multiplication, which forces this form.
Our cost--KL conversions instead involve functions such as
\(\sqrt\kappa\) and \(1-e^{-\kappa}\). The shared question is how tightly
one measurement determines another. Equivalence of finite-dimensional
norms also gives a positive lower endpoint when the fixed input norm is
positive. The examples below show that the cost discrepancy can instead
be zero at positive KL. \Cref{app:moduli-analysis}
explains the terminology's connection to continuity and approximation.

We can ask the same questions for other probability metrics, such as
TV: which cost discrepancies are compatible with a fixed distance,
and what are the sharp conversion functions? The next section will
examine TV alongside KL.

\section{Examples of task cost control at fixed KL}
\label{sec:binary}
We now use point-mass data distributions to illustrate how tightly
fixing forward KL controls task cost. Upper bounds alone leave open
whether the smallest and largest compatible costs are close; here we
characterize the entire interval and show that they need not be.
For small \(\kappa\), its width is approximately \(\kappa\) times the
off-data cost span, so whether it is large depends on the attainable
\(\kappa\), discussed in the next section.

\subsection{The case of point-mass data distributions}
\label{sec:fixed-kl-interval}

We first calculate the full interval of expected costs at a fixed
forward KL for data concentrated on a single successful outcome.
Write \(P=\delta_{x^\star}\), where \(c(x^\star)=0\).
The KL constraint fixes \(Q(x^\star)\); we can choose how to
distribute the remaining probability.

\begin{proposition}[Cost interval for a point-mass reference]
\label{prop:point-mass-cost-interval}
Let \(\cX\) have at least two elements, fix \(x^\star\in\cX\), and
let \(c:\cX\to\R\) satisfy \(c(x^\star)=0\).
For \(P=\delta_{x^\star}\) and \(\kappa\ge0\),
\begin{equation}
 \{\E_Q[c]:\KL(P, Q)=\kappa\}
 =(1-e^{-\kappa})
 \left[\min_{x\ne x^\star}c(x),\,
       \max_{x\ne x^\star}c(x)\right].
 \label{eq:point-mass-cost-interval}
\end{equation}
The endpoints are attained by placing all probability outside
\(x^\star\) on a minimizing or maximizing outcome, respectively.
\end{proposition}

\begin{proof}
By algebra, the KL constraint is \(-\log Q(x^\star)=\kappa\). Thus every
admissible law has the form
\begin{equation}
 Q=e^{-\kappa}\delta_{x^\star}+(1-e^{-\kappa})R,
 \qquad R(x^\star)=0.
 \label{eq:point-mass-mixture}
\end{equation}
Since $c(x^\star)=0$, we have
$\E_Q[c]=(1-e^{-\kappa})\E_R[c]$,
where we are free to choose $R$ subject to $R(x^\star)=0$.
Point masses \(R\) at a cost minimizer or maximizer give the two
endpoints. Mixtures of these point masses give every intermediate value.
\end{proof}

In words, the result states that
KL fixes the mass \(1-e^{-\kappa}\) assigned away from the demonstrated
outcome; its average cost can be anywhere between the smallest and
largest costs available there.

We note in passing that total variation records the same mass:
\[
 \tv{P-Q}=1-Q(x^\star)=1-e^{-\kappa}.
\]
Thus all models in the interval also have the same TV.
The width of the cost interval tells us exactly how much task performance
can vary while both distributional discrepancies remain unchanged.

\subsection{Sequence examples: incorrect steps and final answers}
\label{sec:binary-intervals}

We now ask how models with the same sequence KL can differ in two
assessments of a generated solution: its number of incorrect steps
and the correctness of its final answer. Applying
\cref{prop:point-mass-cost-interval} to each cost gives the exact
range of possible performance.

Now take \(\cX=\{0,1\}^T\), with \(T\ge2\), and set
\(x^\star=0^T\).
Think of this as a simplified model of generating a solution to a
mathematical problem.
For example, let a one record an error in a $T$-step argument, and let
the final bit record whether the derived answer is incorrect.
With this, there are two natural assessments of a derivation $x_{1:T}\in \{0,1\}^T$:
\begin{equation}
 c_{\mathrm{cnt}}(x_{1:T})=\sum_{t=1}^T x_t,\qquad
 c_{\mathrm{fin}}(x_{1:T})=x_T.
 \label{eq:binary-cost}
\end{equation}
The \emph{error count} \(c_{\mathrm{cnt}}\) measures how many steps are wrong. The
\emph{terminal cost} \(c_{\mathrm{fin}}\) measures whether the final
answer is wrong; earlier mistakes matter only through their effect
on that answer.
Under $P=\delta_{0^T}$, the data distribution has zero cost under
either assessment.

\begin{proposition}[Error count at fixed KL]
\label{prop:binary-moduli}
For \(P=\delta_{0^T}\) and every \(\kappa\ge0\),
\begin{equation}
 \{\E_Q[c_{\mathrm{cnt}}]:\KL(P, Q)=\kappa\}
 =[1-e^{-\kappa},\,T(1-e^{-\kappa})],
 \label{eq:binary-kl-level}
\end{equation}
where \(Q\) ranges over all laws on \(\{0,1\}^T\).
\end{proposition}

\begin{proof}
Outside \(0^T\), the smallest count is \(1\), attained at (say)
\(10^{T-1}\), and the largest is \(T\), attained at \(1^T\).
Apply \cref{prop:point-mass-cost-interval}.
\end{proof}

At any positive KL, the largest possible expected count is \(T\)
times the smallest. Even a sharp upper bound leaves this entire
factor-\(T\) spread. A model that makes just one mistake on each
erroneous sequence and a model that makes \(T\) mistakes can have
exactly the same likelihood.

For terminal cost, an erroneous sequence can still end with a correct
answer. This changes the lower endpoint to zero.

\begin{proposition}[Terminal error at fixed KL]
\label{prop:terminal-moduli}
For \(P=\delta_{0^T}\), \(T\ge2\), and every \(\kappa\ge0\),
\begin{equation}
 \{\E_Q[c_{\mathrm{fin}}]:\KL(P, Q)=\kappa\}
 =[0,\,1-e^{-\kappa}],
 \label{eq:terminal-kl-level}
\end{equation}
where \(Q\) ranges over all laws on \(\{0,1\}^T\).
\end{proposition}

\begin{proof}
Outside \(0^T\), terminal cost has minimum \(0\), attained at
\(10^{T-1}\), and maximum \(1\), attained at \(1^T\).
Apply \cref{prop:point-mass-cost-interval}.
\end{proof}

At the same KL, one generator can always give the correct final answer,
while another has error probability \(1-e^{-\kappa}\). Here the cost
range is fixed at \([0,1]\), independently of the horizon.

The interval widths are governed by $1-e^{-\kappa}$, which approaches
one for large $\kappa$. For small \(\kappa\),
\[
 1-e^{-\kappa}=\kappa+O(\kappa^2)\sim\kappa.
\]
The expected count therefore ranges from approximately \(\kappa\)
to \(T\kappa\), and final-answer error ranges from zero to approximately
\(\kappa\). Small sequence KL gives small terminal error, but models
with that same KL can differ throughout the permitted interval.
At \(\kappa=1\), that interval extends from zero to approximately
\(0.632\). Whether such a guarantee is useful depends on the task's
required accuracy.

\subsection{Attaining the cost bounds: recovery or persistence}
\label{sec:endpoint-generators}
We now construct two sequence laws that attain the lower and upper cost
bounds in the preceding subsection. The construction in the proof of
\cref{prop:point-mass-cost-interval} suggests taking
\begin{equation}
 \begin{split}
 Q_-&=e^{-\kappa}\delta_{0^T}
       +(1-e^{-\kappa})\delta_{10^{T-1}},\\
 Q_+&=e^{-\kappa}\delta_{0^T}
       +(1-e^{-\kappa})\delta_{1^T}.
 \end{split}
 \label{eq:endpoint-generators}
\end{equation}
The law $Q_{-}$ attains the lowest cost, while $Q_{+}$ attains the
highest cost for both cost functions considered in the preceding
subsection, among all laws $Q$ with
$\KL(\delta_{0^T}, Q)=\kappa$.

Generation under these laws has a similar form.
Both start with \(X_1\sim\operatorname{Bernoulli}(1-e^{-\kappa})\).
Under \(Q_-\), one continues with \(X_{t+1}=0\); under \(Q_+\),
one continues with \(X_{t+1}=X_t\).

Under both laws, only the last bit is needed to generate the next one.
This is the Markov property:

\begin{definition}[Markov process]
A finite-state discrete-time \emph{stochastic process}
is a collection of random variables $(X_t)_{t=1}^T$ taking values in a
finite set $\cX$. It has the \emph{Markov property}, and is called a
\emph{Markov process}, if for every $t<T$ there is a map
$K_t:\cX^2\to[0,1]$ satisfying
$\sum_{y\in\cX}K_t(x,y)=1$ for every $x\in\cX$, such that, for every
$y\in\cX$, $\P$-almost surely,
\[
  \P(X_{t+1}=y\mid X_1,\ldots,X_t)=K_t(X_t,y)\,.
\]
The rule \(K_t\) is called the process's \emph{transition kernel}.
If the same kernel \(K\) is used at every time, the process is
\emph{time-homogeneous}.
\end{definition}
After identifying $\cX$ with $[d]$, we can represent $K_t$ by a
$d\times d$ nonnegative matrix whose $(i,j)$th entry is $K_t(i,j)$
and whose rows sum to one.

Both laws can be generated by time-homogeneous Markov processes. The
following table gives one choice of transition kernel for each law,
specified by $Q(X_{t+1}=1\mid X_t=x)$, along with the resulting costs:
\begin{center}
\begin{tabular}{lcc}
\toprule
 & \(Q_-\) & \(Q_+\)\\
\midrule
\(Q(X_{t+1}=1\mid X_t=x)\) & \(0\) & \(x\)\\
Expected error count & \(1-e^{-\kappa}\) & \(T(1-e^{-\kappa})\)\\
Final-answer error & \(0\) & \(1-e^{-\kappa}\)\\
\bottomrule
\end{tabular}
\end{center}
We can interpret the generation law of $Q_-$ as recovering from an
initial error, and that of $Q_+$ as persisting that error.

By the preceding results, the intervals also determine both moduli of
continuity introduced in
\cref{sec:discrepancy-comparison}. \Cref{ex:discrepancy-comparison}
asks for these conversions and their interpretation: small expected
error count forces small KL here, while perfect final answers permit
arbitrarily large KL.
In addition, \cref{ex:binary-count} compares the
Pinsker and exact-TV bounds in a variant with an error opportunity
at every step; the additional calculations make the same distinction
between sharpening a guarantee and choosing a better generator.

\section{What does likelihood fitting guarantee?}
\label{sec:growing-horizon}

What does likelihood fitting guarantee about the generated solutions
when data and computation are limited? We first ask how sequence length
and parameter sharing affect learning accuracy, and why lower per-token
loss need not mean better task performance. We then examine the accuracy
required of the optimizer and how optimization error combines with
finite-data error. Together, these comparisons identify what likelihood
fitting controls and which differences between generators it leaves
unresolved.

\subsection{Repeated parameters and new parameters}
\label{sec:parameter-sharing}

Increasing the length of the training sequences gives us more
observations, but fitting longer sequences requires predicting more
tokens. Does the additional information compensate for the additional
predictions? We compare cases where new positions reuse unknown
parameters with cases where they introduce additional unknown parameters,
then return to the resulting task-cost ranges.

Fix a data law \(P\) and a model class \(\mathcal Q\) over infinite
token sequences. Write \(P_{1:T},Q_{1:T}\) for their length-\(T\)
prefix laws. At each length \(T\), a learner receives \(n\) independent
training sequences \(X^{(1)}_{1:T},\ldots,X^{(n)}_{1:T}\) and fits
a model \(\widehat Q^{(n,T)}\), a law over infinite token sequences.
Define the \emph{expected sequence KL for $n$ training sequences of
length $T$} as
\begin{equation}
 \kappa_n(T)=\E\!\left[
 \KL(P_{1:T},\widehat Q_{1:T}^{(n,T)})\right].
 \label{eq:expected-sequence-kl}
\end{equation}
The expectation averages over the training sample with law
\(P_{1:T}^{\otimes n}\) and any learner randomness.
The fitted model can change with both \(n\) and \(T\).
We suppress the superscript \((n,T)\) when the training sample size
and length are clear, writing \(\widehat Q=\widehat Q^{(n,T)}\).
Subscripts such as \(1:T\) indicate the prefix being evaluated.
Unless specified otherwise, we keep \(n\) fixed as \(T\) increases.

For a fixed prediction task, additional training observations cannot
worsen the best attainable statistical accuracy: a learner can ignore
them. In particular, from \(n\) longer sequences it can keep only the
shorter prefixes and run its original procedure. Here we also enlarge
the prediction task. For an unchanged model \(Q\), the KL chain rule
in \cref{eq:population-kl} gives
\[
 \KL(P_{1:T}, Q_{1:T})
 \le \KL(P_{1:T+1}, Q_{1:T+1}).
\]
Learning from the additional observations can improve the fitted
model. We ask whether that improvement compensates for the new
prediction terms.

For models fitted by iterative optimization, exploiting better
statistical accuracy also calls for a sufficiently small optimization
error. Thus the appropriate fitting tolerance depends on both the
amount of data and the sequence length.
\Cref{sec:finite-class-fitting} makes this balance explicit.

A simple estimation problem exhibits both possibilities. First suppose
all tokens are independent \(\operatorname{Bernoulli}(p)\) variables.
This defines $P$.
Estimate their common parameter using the smoothed frequency
\[
 \widehat p=
 \frac{1+\sum_{i=1}^n\sum_{t=1}^T X_t^{(i)}}{nT+2}.
\]
(Smoothing avoids the infinite expected KL of the
ordinary empirical frequency; see \cref{ex:sharing}.)
The fitted model $\widehat Q^{(n,T)}$ is
the independent Bernoulli bit model with parameter \(\widehat p\).
For fixed \(p\in(0,1)\), the large-sample calculation gives
\begin{equation}
 \E[(\widehat p-p)^2]\sim\frac{p(1-p)}{nT},
 \qquad
 \kappa_n(T)=T\,\E[\kl(p,\widehat p)]
 \longrightarrow\frac1{2n}
 \quad(T\to\infty).
 \label{eq:shared-bernoulli}
\end{equation}
\cref{ex:sharing} asks you to prove this and gives hints on the calculation.
Note that each new
position supplies another observation of the same unknown parameter.
The gain in estimation accuracy compensates for the longer prediction
task, giving $\lim_{T\to\infty}\kappa_n(T)=1/(2n)$.

At the other extreme, let the tokens be independent with unrelated
parameters \(p_t\). In particular, for $t\ge 1$,
$X_t$ is Bernoulli with parameter $p_t$.
Estimate each \(p_t\) separately from the \(n\)
training sequences, using
\[
 \widehat p_t=\frac{1+\sum_{i=1}^n X_t^{(i)}}{n+2}.
\]
Each parameter still receives only \(n\) observations. For fixed \(n\),
each estimation risk remains positive; if \(p_t\in[a,1-a]\) for a
fixed \(0<a<1/2\), their sum grows linearly with \(T\).
The large-\(n\) calculation identifies the leading coefficient:
for fixed \(T\) and \(p_1,\ldots,p_T\in(0,1)\),
\begin{equation}
 \kappa_n(T)=\sum_{t=1}^T\E[\kl(p_t,\widehat p_t)]
 \sim\frac{T}{2n}
 \quad(n\to\infty).
 \label{eq:unshared-bernoulli}
\end{equation}
Thus the leading estimation cost grows with the number of positions.
The calculation describes the same fitting method in two model
classes: one shared parameter versus one parameter per position.

\paragraph{Partial sharing.}
Fix a known partition \(G_1,G_2,\ldots\) of the positive integer
positions, ordered by first occurrence. Tokens are independent, with
\(\P(X_t=1)=p_j\) for \(t\in G_j\).
Let \(d(T)\) be the number of groups represented among the first \(T\)
positions, and let \(m_j=|G_j\cap[T]|\). Pool all observations belonging
to the same group and fit independent bits with the estimated probabilities
\begin{equation}
 \widehat p_j=
 \frac{1+\sum_{i=1}^n\sum_{t\in G_j\cap[T]}X_t^{(i)}}{nm_j+2}.
 \label{eq:pooled-frequency}
\end{equation}
The resulting sequence KL grows with the number of unknown parameters:
each contributes approximately \(1/(2n)\).

\begin{restatable}[Sequence KL under parameter sharing]{proposition}{sharingrisk}
\label{prop:sharing-kl}
In the grouped independent-bit model, fit each probability by the pooled
smoothed frequency from \(n\) independent length-\(T\) training sequences.
Fix \(0<a<1/2\) and suppose \(p_j\in[a,1-a]\). Then
\begin{equation}
 \kappa_n(T)\asymp_a\frac{d(T)}n,\qquad n,T\ge1.
 \label{eq:sharing-rate}
\end{equation}
Here \(\asymp_a\) means upper and lower bounds by positive constant
multiples, with constants depending only on \(a\).
If \(n\min_{1\le j\le d(T)}m_j\to\infty\), then more precisely
\begin{equation}
 \kappa_n(T)\sim\frac{d(T)}{2n}.
 \label{eq:sharing-leading}
\end{equation}
\end{restatable}

Each independently estimated parameter therefore contributes approximately
\(1/(2n)\), regardless of how many positions use it.
To see the cancellation, group \(j\) supplies \(nm_j\) training
observations and contributes \(m_j\) terms to the test-sequence KL:
\[
 \kappa_n(T)=\sum_{j=1}^{d(T)}m_j\,\E[\kl(p_j,\widehat p_j)]
 \approx\sum_{j=1}^{d(T)}m_j\frac1{2nm_j}
 =\frac{d(T)}{2n}.
\]
The condition in the proposition says that every pooled sample size
becomes large, so this approximation is accurate across all groups.
\Cref{ex:sharing} develops the proof;
\cref{app:growing-horizon} gives the proof.

The two extremes are \(d(T)=1\) and \(d(T)=T\).
Intermediate sharing patterns give intermediate growth of the
estimation cost. For example, using one new parameter for each
successive block of lengths \(1,2,3,\ldots\) gives
\(d(T)\) of order \(\sqrt T\).
A fixed finite collection of known candidate models offers another
possibility: exact identification can become sufficiently reliable
that \(\kappa_n(T)\) tends to zero (even with $n$ fixed).
\Cref{ex:bernoulli-identification} asks for this calculation with
two Bernoulli candidates.

These comparisons keep the number of independent sequences \(n\)
fixed, allowing the training-token budget \(nT\) to grow.
If instead we fix \(L=nT\), write
\(\kappa'_L(T):=\kappa_{L/T}(T)\). In the fully shared model,
substituting \(n=L/T\) and \(d(T)=1\) into \cref{eq:sharing-leading}
gives
\begin{equation}
 \kappa'_L(T)\sim\frac{1}{2(L/T)}=\frac{T}{2L}
 \qquad(L\to\infty).
 \label{eq:fixed-token-sharing}
\end{equation}
Thus, at a fixed large training-token budget, increasing \(T\)
worsens expected sequence KL, which grows approximately linearly.
\Cref{ex:sharing-horizon} also asks for the unshared comparison.

For the class with \(p_j\in[a,1-a]\), the order \(d(T)/n\) is also
\emph{minimax optimal}: no learner improves its worst-case dependence
on \(n\) and \(d(T)\); see the \hyperref[note:sharing-minimax]{bibliographical notes}.

\begin{restatable}[Learning rates and task-cost intervals]{proposition}{sharingfullcost}
\label{prop:sharing-full-cost}
For any \(d\) with \(1\le d(T)\le T\), there exist a data law,
terminal and counting costs, and estimator families such that, for all
\(n\ge1,T\ge2\), the estimators have common expected sequence KL
\(\kappa_n(T)\asymp d(T)/n\) and their expected costs range over exactly
\[
 \left[0,1-e^{-\kappa_n(T)/2}\right]
 \quad\text{and}\quad
 \left[1-e^{-\kappa_n(T)/2},T(1-e^{-\kappa_n(T)/2})\right],
\]
respectively.
\end{restatable}
To get this result,
combine the grouped Bernoulli construction with \cref{sec:binary},
using tokens \(X_t=(U_t,V_t)\), with costs computed from \(V_{1:T}\).
\Cref{ex:sharing-costs} develops the proof.

\paragraph{What should we expect for transformers?}
In a mathematical solution, many steps apply familiar rules, while
later steps may require combining definitions and intermediate results
across a long prefix. The statistical question is how much useful shared
structure training can discover. The computational question is whether
the architecture can carry out the required combinations with its
available depth, width, and memory.
The Bernoulli examples isolate the statistical benefit when the sharing
pattern is known. \Cref{ex:copying-context} studies a separate
obstruction: a short copying rule whose required memory grows.

For mathematical reasoning, a plausible working hypothesis is
substantial reuse together with additional demands as problems grow.
The grouped model gives one way to quantify that hypothesis; the
number \(d(T)\) counts its independently estimated probabilities.
Which sharing is useful and executable in a transformer is a question
about the task and the learned computation.
\Cref{ex:positionwise-loss} studies a missed parity dependency as an
exercise in model-class restrictions and optimization accuracy.
Lecture~18 develops the statistical horizon comparison further.

\subsection{The pitfalls of reporting per-token log loss}
\label{sec:late-predictions}

As sequence length grows, log loss per token can vanish while
final-answer error stays constant. For the all-zero law \(P\) and
persistent generator \(Q_+\) from \cref{eq:endpoint-generators},
with fixed \(\kappa>0\),
\begin{equation}
 \frac{\KL(P_{1:T},(Q_+)_{1:T})}{T}
 =\frac{\kappa}{T}\longrightarrow0,
 \qquad
 \E_{(Q_+)_{1:T}}[c_{\mathrm{fin}}]=1-e^{-\kappa}.
 \label{eq:endpoint-normalization}
\end{equation}
The lower per-token loss comes entirely from dividing the same
sequence loss by a larger \(T\).
When comparing different sequence lengths, lower per-token log loss
alone is therefore insufficient evidence of better final-answer accuracy.

\subsection{Likelihood fitting at finite accuracy}
\label{sec:likelihood-pitfall}

We use loss per predicted token as the optimization objective.
For \(n\) observed sequences, define empirical \emph{per-token}
log loss
\[
 \widehat L(Q)=\frac1{nT}\sum_{i=1}^n\sum_{t=1}^T
 -\log Q(X_t^{(i)}\mid X_{<t}^{(i)}).
\]
This averages the losses of the \(nT\) token predictions.
We use this normalization despite the warning in
\cref{sec:late-predictions}: per-token loss can improve as sequence
length grows without task performance improving.
There are two reasons for using it. First, when token losses have
comparable scales, normalization removes the explicit factor of \(T\)
that would otherwise multiply the objective and its gradient.
For gradient descent, this avoids a corresponding adjustment of the
stepsize and helps us use common optimization settings across sequence
lengths. Second, log loss is conventionally reported per token.\endnote{\textbf{Unequal response lengths.}
For response \(y_i\) to prompt \(u_i\), let \(T_i\) count its predicted
tokens and put \(\ell_i(Q)=-\log Q(y_i\mid u_i)\).
Two averaging conventions are
\[
 \widehat L_{\rm token}(Q)=\frac{\sum_i\ell_i(Q)}{\sum_iT_i},
 \qquad
 \widehat L_{\rm response}(Q)=\frac1n\sum_i\frac{\ell_i(Q)}{T_i}.
\]
The first weights predicted tokens equally; the second weights each
response's average loss equally. Prompt and padding tokens excluded from
the loss are also excluded from \(T_i\). Both conventions reduce to the
formula in \cref{sec:likelihood-pitfall} when every \(T_i=T\).
Both are used in supervised fine-tuning; see the
\hyperref[note:sft-normalization]{bibliographical notes}.}

When an optimizer is run with finite compute, it generally does not
return a model with the lowest possible empirical loss.
The \emph{optimization error} of its returned model \(\widehat Q\) is
\begin{equation}
 E(\widehat Q):=\widehat L(\widehat Q)
 -\inf_{Q\in\mathcal Q}\widehat L(Q).
 \label{eq:approximate-fit-error}
\end{equation}
We say that \(\widehat Q\) satisfies the \emph{optimization tolerance}
\(\eta\ge0\) if \(E(\widehat Q)\le\eta\), equivalently,
\begin{equation}
 \widehat L(\widehat Q)\le\inf_{Q\in\mathcal Q}\widehat L(Q)+\eta.
 \label{eq:approximate-fit}
\end{equation}
Establishing whether a returned model satisfies this tolerance is a
\emph{certification question}. For example, suppose
\(\mathcal Q=\{Q_\theta:\theta\in\R^d\}\) and
\(f(\theta)=\widehat L(Q_\theta)\) is differentiable and strongly convex
with known parameter \(\mu>0\). Then
\[
 E(Q_\theta)\le\frac{\|\nabla f(\theta)\|_2^2}{2\mu},
\]
so \(\|\nabla f(\theta)\|_2\le\sqrt{2\mu\eta}\) certifies the tolerance
(\cref{ex:optimization-certificate}). A small gradient alone gives no
such guarantee for a general nonconvex loss.
We still wanto to study the consequences of positive optimization error
even when its size is unknown to the optimizer because this helps with
understanding the impact of putting more effort into optimization.

Because of the normalization by \(T\), a per-token tolerance \(\eta\)
only guarantees a per-sequence tolerance \(T\eta\).
Our earlier cost bounds motivate controlling the per-sequence error.
Maintaining a fixed per-sequence tolerance as \(T\) increases requires
increasingly accurate per-token optimization, which may require more
compute. Whether a fixed number of passes through the dataset suffices
depends on the optimization problem and how the number of updates
scales with \(T\).\endnote{\textbf{Fixed passes and finer accuracy.}
Consider smooth, strongly convex per-token objectives with constants
uniform in \(n,T\), bounded initial gaps, and unbiased stochastic gradients
with uniformly bounded variance. With suitable decreasing stepsizes,
stochastic gradient descent gives expected optimization error
\(O(1/K)\) after \(K\) updates. If batches contain a fixed number of
uniformly sampled token losses, \(K=\Theta(nT)\) updates process a fixed
multiple of the dataset's token count and give expected per-sequence
error \(O(1/n)\). This counts token-loss evaluations; their computational
cost can depend on context length. The bound averages over optimizer
randomness and does not certify an individual run. Full-gradient descent uses one pass per
update, so fixed passes give fixed \(K\). Even for
\(f(\theta)=\theta^2/2\), starting at \(\theta_0=1\) with stepsize
\(1/2\) gives error \(2^{-2K-1}\), independent of \(T\); reaching
\(s/T\) for fixed \(s>0\) requires order \(\log T\) updates.
Thus passes alone do not determine the achieved tolerance. The
stochastic rate also need not hold for a general neural-network loss.
Lecture~10 returns to when local gradients ensure global convergence;
see the \hyperref[note:optimization-accuracy]{bibliographical notes}
for the assumptions and rate.}\label{en:optimization-compute}

Positive optimization error can translate into poor task performance,
just as sequence-level KL error can. Since \(\eta\) is measured per
token, the relevant sequence-level quantity is \(T\eta\).
The simplest illustration uses the deterministic data law
\(P=\delta_{0^T}\) from \cref{sec:endpoint-generators}: empirical and
population log loss coincide, so there is no statistical error to
account for. Reusing that example gives the following result.

\begin{proposition}[Different task costs at the same optimization tolerance]
\label{prop:endpoint-optimization}
Let \(T\ge2\). Take the laws \(Q_-,Q_+\) in \cref{eq:endpoint-generators} with
\(\kappa=T\eta\), and let the model class contain \(P,Q_-,Q_+\).
For every sample of \(n\) sequences from \(P=\delta_{0^T}\),
\begin{equation}
 \begin{gathered}
 \widehat L(P)=0,\qquad
 \widehat L(Q_-)=\widehat L(Q_+)=\eta,\\
 \E_{Q_-}[c_{\mathrm{fin}}]=0,\qquad
 \E_{Q_+}[c_{\mathrm{fin}}]=1-e^{-T\eta}.
 \end{gathered}
 \label{eq:endpoint-optimization}
\end{equation}
Both generators satisfy \cref{eq:approximate-fit}.
\end{proposition}

\begin{proof}
Each data sequence is \(0^T\). Both generators assign it probability
\(e^{-T\eta}\), so their loss per token is \(\eta\).
Their terminal costs follow from \cref{eq:endpoint-generators}.
\end{proof}

For small \(T\eta\), the permitted final-answer error is approximately
\(T\eta\). For example, \(\eta=10^{-3}\) and \(T=1000\) permit
anything from zero terminal error to about \(0.632\).
More demonstrations of \(0^T\) leave this choice unchanged.

If the class consists of \(Q_-,Q_+\) at a fixed \(\kappa>0\), then
the data law \(P\) is outside the class. Both generators minimize
empirical and population log loss, yet their terminal costs are zero
and \(1-e^{-\kappa}\). Thus, in addition to optimization error,
\emph{approximation error} can also leave a
consequential choice between equally fitted models.
\Cref{ex:fitting-support} develops this comparison with repeated
errors, and \cref{ex:positionwise-loss} studies a missed dependency
in a parity task.

\subsection{What does finite-sample likelihood fitting guarantee?\suppmark}
\label{sec:finite-class-fitting}

Likelihood fitting gives a positive statistical guarantee
on task performance.
The following stylized result bounds the task cost of choosing a model
with nearly minimal empirical log loss from a finite class.
Allowing optimization slack, as in the previous subsection, lets us ask
how optimization effort should depend on sample size.

\begin{restatable}[Approximate maximum likelihood controls cost]{theorem}{finitemle}
\label{thm:finite-mle}
Let \(\mathcal Q\) be a finite class containing \(P\), and observe
\(n\) independent sequences with law \(P\).
Let \(0\le c\le C_{\max}\).
Suppose \(\widehat Q\in\mathcal Q\) satisfies
\cref{eq:approximate-fit} for \(\eta\ge0\).
For \(0<\delta<1\), put
\[
 b=\frac{2\log(|\mathcal Q|/\delta)}n+T\eta.
\]
With probability at least \(1-\delta\),
\begin{equation}
 d_c(P,\widehat Q)\le2\sigma_P(c)\sqrt b+C_{\max}b.
 \label{eq:finite-mle}
\end{equation}
In particular, if \(\E_P[c]=0\), then
\(\E_{\widehat Q}[c]\le C_{\max}b\).
\end{restatable}

The proof in \cref{app:finite-mle} uses the same Hellinger-to-cost
inequality as \cref{prop:forward-variance}
(\cref{prop:hellinger-cost}). The earlier result bounds squared
Hellinger distance by forward KL; here the likelihood argument bounds
it directly by \(b\). Forward KL itself need not be finite, so the
earlier proposition alone does not give this theorem.

For zero-cost data and \(0\le c\le1\), the theorem separates the two
contributions to its guarantee:
\begin{equation}
 \E_{\widehat Q}[c]\le
 \underbrace{\frac{2\log(|\mathcal Q|/\delta)}n}_{\text{statistical contribution}}
 +\underbrace{T\eta}_{\text{optimization contribution}}.
 \label{eq:finite-mle-zero-cost}
\end{equation}
At fixed model class and confidence level, the statistical contribution
is of order \(1/n\) and remains in the bound even with exact optimization.
If optimization effort is insufficient, the optimization contribution
\(T\eta\) may dominate. To keep it no larger than the statistical
contribution, it suffices to achieve
\begin{equation}
 \eta\le\frac{2\log(|\mathcal Q|/\delta)}{nT}.
 \label{eq:statistical-accuracy}
\end{equation}
Substitution into \cref{eq:finite-mle-zero-cost} gives
\(\E_{\widehat Q}[c]\le4\log(|\mathcal Q|/\delta)/n\), with probability
at least \(1-\delta\).
\emph{Optimizing to statistical accuracy} means keeping
the optimization contribution on the scale of that statistical term.
Even at fixed sequence length, retaining this statistical rate calls
for a decreasing per-token tolerance as the number of examples grows.

The example in \cref{prop:endpoint-optimization} shows that the
optimization contribution has the correct order when \(T\eta\) is small.

Lecture~18 develops this comparison for imitation learning.
The comparison in \cref{sec:parameter-sharing} fixes infinite-sequence
models and varies the length of the observed prefixes. Its Bernoulli
examples isolate parameter sharing; the present finite-class theorem
also accounts for optimization tolerance. Fixing the number of sequences
gives a different comparison from fixing the token budget.

\section{Beyond likelihood fitting can we improve task performance?}
\label{sec:better-control}
The main conclusion of the previous section is that
in terms of task performance, likelihood fitting
may leave a large room for improvement.
In this section we investigate how to decrease this room.
For this, we consider two broad approaches.
The first considers changing the objective to incorporate information
about the consequences of generation,
while the second one gives the models
more structure to support
recovery from errors. Recovery then raises a further question: which
information should be forgotten, and which must be preserved?

\subsection{Evaluate consequences of generation}
\label{sec:generation-objectives}

\paragraph{Direct cost feedback.}
In some cases, the cost function is available for evaluation.
That is, there is a cheap way to get the value of $c$ at any sequence $x_{1:T}$.
(Note that this is not always possible.
This is the case, for example, when the value of $c(x_{1:T})$ involves asking a human, and no human
is available to perform this task.)
However, when $c$ is available for measurement, this creates a new opportunity
to get a good model. One can simple run an optimization process
that aims to minimize
\begin{equation}
 \min_\theta\E_{Q_\theta}[c].
 \label{eq:direct-cost}
\end{equation}
The procedures that aim to perform this optimization are discussed later. In particular,
Lecture~19 develops reinforcement learning with this feedback, including the policy-gradient
method for optimizing such an objective. Lectures~20--21 study preference
feedback and optimization with learned evaluators.

Note that in practice, as alluded to beforehand, one often combines the various data sources
and objectives. This should make perfect sense when all the objectives are aligned.
This would be case when teacher forcing with positive data is used alongside with the above
optimization objective. The reason to combine the objectives is to take advantage
of their complementary strengths.
An example of how access to cost can help consider the problem of minimizing the expected terminal
cost $c_{\rm fin}$. As we have seen, $Q_+$ and $Q_-$ of \cref{sec:endpoint-generators} both have the same
log-likelihood, while the expected cost of $Q_-$ is zero, and the expected cost of $Q_+$ is \(1-e^{-\kappa}\).
If the expected cost can be calculated up to sufficient accuracy (less than the performance gap
between the models), the better, $Q_-$ can be chosen, which is a large win compared to not using cost information.
TODO: Are there complementary strengths? I would think so, but we better develop some examples. What is nontrivial is why teacher forcing with positive data can help if we can also optimize the cost directly. The answer may be that the other objective has better optimization properties. Also, there is the ``why not'' answer. Two things are better than one (we are assuming they are compatible). Anyhow, need to add some exercise about this.


\paragraph{Predicting further ahead.}
\label{sec:multistep}

Positive demonstrations can supervise predictions several steps into the
future as well as predictions of the next token. The model must predict
\(X_{t+k}\) from \(X_{\le t}\), before the intervening tokens are supplied.
For an observed sequence \(x_{1:T}\) and a model with \(Q(x_{1:T})>0\),
define the \emph{\(k\)-step prediction loss}, for \(1\le k\le T\), by
\begin{equation}
  \cL_k(Q)
  =-\sum_{t=0}^{T-k}
  \log Q(X_{t+k}=x_{t+k}\mid X_{\le t}=x_{\le t}).
  \label{eq:multistep-loss}
\end{equation}
The prefix is empty when \(t=0\). In particular,
\(\cL_1(Q)=-\log Q(x_{1:T})\) is the usual sequence log loss.
One can add \(\lambda\cL_k(Q)\), with \(\lambda>0\) and \(k\ge2\),
to this objective and average over the demonstrations.

For the laws in \cref{sec:endpoint-generators}, evaluated on \(0^T\),
\begin{equation}
 \cL_2(Q_-)=0,\qquad \cL_2(Q_+)=\kappa.
 \label{eq:endpoint-multistep}
\end{equation}
Their next-token losses agree. At \(\kappa>0\), any positive weight
on \(\cL_2\) therefore favors \(Q_-\), which recovers from its initial
error, over \(Q_+\), which preserves it. The calculation is in
\cref{ex:multistep-benefit}.\endnote{\textbf{A separate benefit under misspecification.}
A finite model class may contain a law with the correct marginal of
\(Y=f(X_{1:T})\), even when no model matches the full sequence law.
The finite-class guarantee in \cref{thm:finite-mle} then applies to
fitting the marginal likelihood of \(Y\) and any cost determined by
\(Y\). For a Markov generator and \(T=mk\), the likelihood of
\(Y=(X_k,X_{2k},\ldots,X_T)\) is the product of the \(k\)-step
predictions at starts \(0,k,\ldots,T-k\). The guarantee therefore applies
to this subsampled likelihood objective.}\label{en:multistep-misspecification}

TODO: Find substantive conditions under which adding a short-lookahead
loss improves task performance over next-token fitting, beyond the
\(Q_-,Q_+\) example. A bound in terms of achieved multistep loss alone
does not establish this comparison.

The computational challenge is to evaluate and optimize these losses.
For a general autoregressive model,
\begin{equation}
\begin{split}
 Q_\theta(X_{t+k}=a\mid x_{\le t})
 =\sum_{y_{1:k-1}\in[N]^{k-1}}
 &\left[\prod_{j=1}^{k-1}
 Q_\theta(y_j\mid x_{\le t},y_{<j})\right]
 \cdot Q_\theta(a\mid x_{\le t},y_{1:k-1}).
\end{split}
\label{eq:multistep-marginal}
\end{equation}
Direct evaluation sums over \(N^{k-1}\) possible intervening sequences;
for a large vocabulary, even \(k=2\) can be costly.
Supplying the observed intervening tokens instead would return to a
next-token prediction problem. An auxiliary prediction head can instead
be trained to predict the future token directly from the prefix, without
evaluating this sum. Its predictions need not agree with the generator's
\(k\)-step marginals.\endnote{\textbf{Auxiliary prediction heads and consistency.}
A separate head \(h_{\theta,k}(a\mid x_{\le t})\) can be trained directly
on observed prefix--target pairs. This is auxiliary supervision, used
in multi-token prediction. To represent the original generator's
\(k\)-step marginal, the head must satisfy
\[
 h_{\theta,k}(a\mid x_{\le t})
 =Q_\theta(X_{t+k}=a\mid x_{\le t}).
\]
Sharing parameters does not enforce this equality. A penalty can encourage
agreement without enforcing it; fitting a head to a fixed teacher's
predictions is a form of \emph{distillation}. For a fixed generator,
future-token samples from its rollouts provide training targets without
enumerating all intervening sequences, at the cost of sampling error.
Exact agreement remains a separate requirement, and accurate auxiliary
heads alone do not guarantee accurate generation. See the
\hyperref[note:multistep-heads]{bibliographical notes} for examples and
references.}\label{en:multistep-heads}

The intervening tokens are latent variables in this prediction problem:
\cref{eq:multistep-marginal} is a conditional likelihood obtained by
summing them out. Lecture~5 returns to the same
operation for hidden states. \Cref{ex:multistep} develops the calculation
and shows how Markov structure makes it efficient.

\paragraph{Fit reverse KL.}
\label{sec:reverse-kl}

Forward KL penalizes failing to cover outcomes that are likely under \(P\).
Reverse KL penalizes generating outcomes to which \(P\) assigns too little
probability. When a model cannot reproduce \(P\), this distinction matters:
we may prefer to generate a smaller variety of successful responses if
covering more responses would also introduce failures.

To quantify the penalty on failures (positive costs), let \(B=\{x:c(x)>0\}\),
\(\varepsilon=P(B)\in(0,1)\), and \(q=Q(B)\). Data processing gives
\begin{equation}
 \KL(Q,P)\ge\operatorname{kl}(q,\varepsilon)
 \ge q\log(1/\varepsilon)-\log 2.
 \label{eq:reverse-failure-prob}
\end{equation}
Thus, reverse KL increases
at least at a linear speed with the failure probability $q=Q(B)$
with steepness $\log(1/\varepsilon)$, which diverges as $\varepsilon\to 0$.
Conversely, a bounded reverse KL forces failures
to become rare when they are rare under \(P\).
\Cref{thm:reverse-cost} gives a related cost guarantee using a subgaussian
variance proxy; \cref{eq:reverse-failure-prob} directly controls failure
probability.

Finite reverse KL requires staying in the support of \(P\). When $P(B)=0$,
finite reverse KL implies zero task cost.
But finite reverse KL can also exclude probability laws that have zero task cost:
\(Q_-\) has zero terminal cost but infinite reverse KL from
\(P=\delta_{0^T}\) (just like \(Q_+\)).
\Cref{ex:kl-directions} develops the failure bound and the complementary
calculations for adding unsuccessful outcomes and omitting successful ones.

\paragraph{The information needed for reverse-KL fitting.}
Forward KL averages conditional discrepancies at data prefixes.
Reversing its arguments averages them at generated prefixes:
\begin{equation}
 \KL(Q, P)
 =\sum_{t=1}^T\E\!\left[
 \KL\bigl(Q(\,\cdot\mid\widetilde X_{<t})
       , P(\,\cdot\mid\widetilde X_{<t})\bigr)\right],
 \qquad \widetilde X_{1:T}\sim Q.
 \label{eq:reverse-chain}
\end{equation}
Thus the generator's own erroneous histories contribute to the objective
\begin{equation}
 \KL(Q_\theta, P)
 =\E[\log Q_\theta(\widetilde X)-\log P(\widetilde X)],
 \qquad \widetilde X\sim Q_\theta.
 \label{eq:reverse-objective}
\end{equation}
For the unknown data law \(P\), the probabilities \(P(\widetilde X)\)
are unavailable. Thus this expression does not directly provide an
objective we can evaluate from the training data.
A \emph{variational representation} avoids requiring these probabilities
and leads to an \emph{adversarial} fitting method.
For \(X\sim P\) and \(Y\sim Q\) on the finite sequence space
\(\mathcal X=[N]^T\),
\begin{equation}
 \KL(Q,P)
 =\sup_{g:\mathcal X\to\R}
 \left\{1+\E[g(Y)]-\E[e^{g(X)}]\right\}.
 \label{eq:reverse-adversarial}
\end{equation}
The function \(g\) is called a \emph{critic}. When \(P,Q\) are positive,
the maximizer is \(g^*(x)=\log(Q(x)/P(x))\).
\Cref{ex:reverse-variational} asks for the proof.
Both expectations can be estimated from samples, without evaluating
\(P(x)\). This suggests alternating updates: fit the critic to maximize
the expression, then fit the generator to minimize it with the critic
fixed.

Restricting the critic to a tractable class yields only a lower bound
on \(\KL(Q,P)\), and fitting it from finite data requires generalization.
Thus a small fitted objective does not certify small reverse KL.
The support requirement remains: if \(Q\) assigns mass where \(P\) is
zero, the unrestricted supremum is infinite; a restricted critic may
fail to detect this.\endnote{\textbf{Learning support from samples.}
The empirical data law gives infinite reverse KL to any generator that
leaves the training set. Exact support recovery can itself require many
demonstrations, even for two candidate rules;
\cref{app:conditional-support} gives the calculation.}

Direct reverse-KL fitting is nevertheless useful for model compression
through \emph{model distillation}: a smaller student approximates an
existing teacher whose probabilities can be evaluated.
Here the target is the teacher's distribution. Both forward and reverse
KL are used for model distillation.\endnote{\textbf{Distillation objectives and training data.}
Forward KL is the conventional soft-target distillation objective:
the student fits the teacher's output probabilities using cross-entropy.
For language models, distinguish the source of prefixes from the source
of response targets. Token-level distillation matches the teacher's
next-token distribution at prefixes from a text or fine-tuning dataset.
Sequence-level forward KL can be fitted by sampling responses from the
teacher and maximizing their likelihood under the student.
Sequence-level reverse KL instead uses student-generated responses and
teacher probability evaluations.

Practical methods modify these schemes. MiniLLM uses instruction-data
prompts, student--teacher mixed generation, supervised fine-tuning for
initialization, and an auxiliary language-modeling loss on a pretraining
corpus. The teacher's original pretraining dataset is not required.
See the \hyperref[note:reverse-fitting]{bibliographical notes} for sources.
}\label{en:distillation}

\subsection{Recovery at unchanged likelihood}
\label{sec:stability}

Suppose the task is to produce a correct final answer at position \(T\).
For a fixed question, let \(c(x)=\mathbf 1\{x\text{ is incorrect}\}\).
The expected cost
\(\E_Q[c(X_T)]\) depends only on the distribution of \(X_T\).
This distribution is the \emph{marginal law at position \(T\)}:
we obtain it from the joint law by summing over the preceding tokens.
More generally,
let \(\mu_t,\nu_t\) denote the marginals under \(P,Q\), respectively:
\[
\mu_t(x)=P(X_t=x),\qquad \nu_t(x)=Q(X_t=x)\,, \qquad t\in [T]\,.
\]
By \cref{prop:tv-cost}, we have
\[
 \left|\E_P[c(X_T)]-\E_Q[c(X_T)]\right|
 \le\tv{\mu_T-\nu_T}.
\]
If the data always end with a correct answer, \(\E_P[c(X_T)]=0\),
so controlling this marginal discrepancy controls the generated
answer's error probability.

An early mistake can be corrected before the final answer. We will
quantify this recovery when the fitted process \(Q\) is Markov,
and then exhibit improved final-answer accuracy at unchanged likelihood.
The same argument controls cumulative costs
\(\sum_{t=1}^T c_t(X_t)\); \cref{ex:cumulative-recovery} asks for the bound.
\Cref{app:history-recovery} extends the result to history-dependent
generators by controlling how long a changed token affects later predictions.

For a law \(\mu\) and kernel \(K\), write
\[
 (\mu K)(y)=\sum_x\mu(x)K(x,y).
\]
Let \(Q\) have transition kernels \(\widehat K_t\), so
\(\nu_{t+1}=\nu_t\widehat K_t\). The data law \(P\) can depend on
the full history. We measure local prediction error at its prefixes,
as in likelihood fitting.

\begin{theorem}[Recovery and final-answer accuracy]
\label{thm:markov-stability}
Suppose that, for \(1\le t<T\), the fitted kernels satisfy
\begin{equation}
 \tv{\alpha\widehat K_t-\beta\widehat K_t}
 \le\rho\tv{\alpha-\beta}
 \quad\text{for all laws }\alpha,\beta,\qquad 0\le\rho\le1,
 \label{eq:markov-contraction}
\end{equation}
and their local errors under \(P\) satisfy
\begin{equation}
 \E_P\!\left[
 \tv{P(\cdot\mid X_{\le t})-\widehat K_t(X_t,\cdot)}
 \right]
 \le\varepsilon_{t+1}.
 \label{eq:markov-local-error}
\end{equation}
Write \(e_t=\tv{\mu_t-\nu_t}\). For every \(c:[N]\to[0,1]\),
\begin{equation}
 \left|\E_P[c(X_T)]-\E_Q[c(X_T)]\right|
 \le e_T
 \le\rho^{T-1}e_1+\sum_{t=2}^T\rho^{T-t}\varepsilon_t.
 \label{eq:stability-recursion}
\end{equation}
If \(\rho<1\), \(e_1\le\varepsilon\), and
\(\varepsilon_t\le\varepsilon\) for all \(2\le t\le T\), then
\begin{equation}
 \max_{1\le t\le T}
 \left|\E_P[c(X_t)]-\E_Q[c(X_t)]\right|
 \le\frac{\varepsilon}{1-\rho}.
 \label{eq:frequency-control}
\end{equation}
\end{theorem}

\begin{proof}
The laws \(\mu_{t+1}\) and \(\mu_t\widehat K_t\) average the true
and fitted next-token conditionals over the same data prefixes:
\[
 \mu_{t+1}=\E_P[P(\cdot\mid X_{\le t})],\qquad
 \mu_t\widehat K_t=\E_P[\widehat K_t(X_t,\cdot)].
\]
Their TV distance is therefore at most \(\varepsilon_{t+1}\).
Insert \(\mu_t\widehat K_t\) between \(\mu_{t+1}\) and
\(\nu_{t+1}=\nu_t\widehat K_t\):
\[
 e_{t+1}
 \le \tv{\mu_{t+1}-\mu_t\widehat K_t}
      +\tv{\mu_t\widehat K_t-\nu_t\widehat K_t}
 \le\varepsilon_{t+1}+\rho e_t.
\]
Iteration and the TV cost bound give \cref{eq:stability-recursion}.
Under the uniform
error bounds, \(e_t\le\varepsilon\sum_{j=0}^{t-1}\rho^j
\le\varepsilon/(1-\rho)\). Apply \cref{prop:tv-cost} to \(c\) at
each time and take the maximum.
\end{proof}

In the final-answer bound, the local error at time \(t\) is multiplied
by \(\rho^{T-t}\). Recovery reduces the effect of
early errors; a mistake in the final prediction has no remaining
steps in which to be corrected. With repeated local errors of size
at most \(\varepsilon\), their combined effect is bounded by
\(\varepsilon/(1-\rho)\), independently of the horizon.
The argument also extends to other distances between distributions,
such as Wasserstein-1 distance with Lipschitz costs and contraction
in that distance.

We can check contraction by comparing the next-token laws at different
current states. For a finite-state kernel, the smallest possible contraction
factor is the \emph{Dobrushin ergodicity coefficient}, given by
\begin{equation}
 \rho(\widehat K)
 =\max_{x,x'}\tv{\widehat K(x,\cdot)-\widehat K(x',\cdot)}.
 \label{eq:dobrushin}
\end{equation}
\Cref{ex:stability} proves this characterization.

\paragraph{Recovery in the binary example.}
We now improve the persistent law \(Q_+\) by giving
each error a chance to be corrected. Here \(c(x)=x\), so
\(c(X_T)=c_{\mathrm{fin}}(X_{1:T})\) is the final-answer error.
The per-position costs \(c_t(x)=x\) give \(c_{\mathrm{cnt}}\).

Fix \(P=\delta_{0^T}\) and \(\kappa\ge0\). Let \(Q_r\) use the same
first-bit distribution as $Q_-$ and $Q_+$, followed by
\[
 \begin{gathered}
 Q_r(X_1=1)=1-e^{-\kappa},\\
 Q_r(X_{t+1}=1\mid X_t=0)=0,\qquad
 Q_r(X_{t+1}=0\mid X_t=1)=r,\quad 1\le t<T,
 \end{gathered}
\]
where \(0\le r\le1\). Once an error is corrected, all later bits
are zero. The choices \(r=0\) and \(r=1\) give \(Q_+\) and \(Q_-\),
respectively.

\begin{proposition}[Recovery improves final-answer accuracy at unchanged likelihood]
\label{prop:recovery}
For this family and \(T\ge2\),
\begin{equation}
 \KL(P, Q_r)=\kappa,\qquad
 \E_{Q_r}[c_{\mathrm{fin}}]
 =(1-e^{-\kappa})(1-r)^{T-1}.
 \label{eq:endpoint-recovery}
\end{equation}
For \(r>0\),
\begin{equation}
 \E_{Q_r}[c_{\mathrm{cnt}}]
 =(1-e^{-\kappa})\frac{1-(1-r)^T}{r}
 \le\frac{1-e^{-\kappa}}r.
 \label{eq:endpoint-recovery-count}
\end{equation}
At \(r=0\), the expected count is \(T(1-e^{-\kappa})\).
\end{proposition}

\begin{proof}
The all-zero path has probability \(e^{-\kappa}\), independently
of \(r\). Also,
\[
 Q_r(X_t=1)=(1-e^{-\kappa})(1-r)^{t-1},
\]
since the first bit must be wrong and each subsequent correction
must fail. Taking \(t=T\) gives terminal error; summing over \(t\)
gives the expected count.
\end{proof}

At a fixed positive recovery probability, terminal error decays
geometrically and expected error count stays bounded as the horizon
grows. An erroneous run has geometric duration with mean \(1/r\)
before truncation at the horizon. Likelihood constrains
the probability of the initial error; the probability \(r\) of returning
to zero determines how likely it is to persist until the final answer.\endnote{\textbf{Which costs
benefit from recovery?} For
\(c_{\mathrm{any}}(x_{1:T})=\mathbf 1\{\exists t:x_t=1\}\),
\(\E_{Q_r}[c_{\mathrm{any}}]=1-e^{-\kappa}\) for every \(r\).
Correcting a mistake improves the final answer and limits the error
count, while leaving unchanged whether a mistake ever occurred.}
\label{en:recovery-cost}

For \(Q_r\), the two rows differ by \(1-r\), so its kernels
after the first token have contraction factor \(1-r\).
The initial discrepancy is \(e_1=1-e^{-\kappa}\), and all subsequent
local errors under \(P\) are zero. The theorem therefore gives
\[
 e_T\le(1-r)^{T-1}(1-e^{-\kappa}),
\]
which is exactly the terminal cost in \cref{eq:endpoint-recovery}.
\Cref{ex:stability} adds fresh errors at every step and computes the
resulting long-run error level. The next question is how to obtain
recovery while retaining information the task needs.

\subsection{What should be forgotten?}
\label{sec:selective-memory}

Random resets give a simple mathematical way to enforce forgetting.
For any kernel \(K\) and reset law \(\pi_{\rm reset}\), the kernel
\[
 K_r(x,\cdot)=(1-r)K(x,\cdot)+r\pi_{\rm reset}(\cdot)
\]
has contraction factor at most \(1-r\).
A reset makes the two initial states use the same distribution.
But forgetting everything can erase exactly the information a task needs.
A proof generator must retain the theorem being proved and useful facts
established earlier.

Limited context and contraction are different properties. The rule
\(X_{t+1}=X_t\) uses only one token of context and preserves its initial
bit forever. During generation, recent tokens can keep transmitting old
information. Restricting a transformer's accessible prefix therefore
does not establish \cref{eq:markov-contraction}.

For prompting, a sliding window creates a more immediate difficulty:
a fact in the original prompt can leave the accessible context before
the model has carried it forward. Subsequent predictions then have no
direct access to it. Generated reminders can preserve information, as
in the copying example, but can also introduce errors into what is retained.
Keeping the original prompt or verified task information separately
provides a more reliable source than repeatedly reconstructing it from
generated text.

A more useful design separates persistent task information \(M\) from
a revisable working state \(R_t\):
\begin{equation}
 S_t=(M,R_t),\qquad
 R_{t+1}=g_M(R_t,\Xi_{t+1}).
 \label{eq:persistent-working-memory}
\end{equation}
Conditional on the same \(M\), the working dynamics can forget
perturbations while \(M\) retains the problem and trusted information.
For example, in proof generation this suggests retaining the problem
statement and checked lemmas while restarting a failed attempt.
The mathematical question is which state differences can be made
contractive without destroying information the task needs.
Lecture~5 develops state-space models as an explicit way to express this separation.


\section{Take-home messages}
\label{sec:take-home}

\begin{enumerate}[leftmargin=*,itemsep=0.3em]
\item Next-token log-loss minimization is sequence maximum likelihood.
Its population excess loss is forward sequence KL.
\item Successful demonstrations specify a distribution to imitate.
The task objective is the expected cost of generated sequences.
TV and KL bounds connect these objectives by limiting the expected
cost discrepancy.
\item At the same forward KL \(\kappa\) and TV, the laws $Q_-$ and $Q_+$
have expected counts \(1-e^{-\kappa}\) and \(T(1-e^{-\kappa})\),
and terminal error probabilities zero and \(1-e^{-\kappa}\).
Approximate likelihood fitting can leave this choice unresolved:
a per-token tolerance \(\eta\) permits per-sequence excess \(T\eta\).
\item Longer training sequences provide more information, while longer
predictions may require learning additional structure. In the grouped
Bernoulli model, \(d(T)\) unknown shared probabilities give expected
sequence KL at most \(d(T)/n\). The persistent law $Q_+$
shows that a vanishing average-loss gap can coexist with unchanged
task error as the horizon grows.
\item Recovery attenuates an initial error; continued prediction errors
contribute according to how long their effects persist. Contraction
bounds both terms. Useful task memory may still need to persist.
\item Forward KL penalizes failing to cover likely data outcomes; reverse
KL penalizes assigning too much probability to outcomes rare under the
target. Adversarial fitting uses samples from both laws, but a small
objective with a restricted critic does not certify small reverse KL.
\item Multi-step objectives integrate over missing intermediate tokens.
Direct cost optimization instead uses evaluations on model generations.
\end{enumerate}

\clearpage
\printendnotes

\section*{Bibliographical notes}
\phantomsection
\label{sec:bibliographical-notes}

The probability and KL chain rules continue Lecture~3.
The cost bounds and binary calculations here follow directly from those
identities. The zero-cost forward-KL refinement makes explicit information
that a general Pinsker conversion discards.

\paragraph{Likelihood, task discrepancy, and imitation learning.}
\Cref{prop:forward-variance} has Bernstein form, with a variance-dependent
square-root term and a range-dependent linear term.
\citet[Theorem~3.1]{foster2024} obtain this form using trajectory-level
Hellinger distance and expert-action variance. Their lower-order term
retains the total-return range \(R\), up to logarithms.
The technical gain in their Section~3.2 is to avoid an extra horizon
factor: summing \(H\) advantages (changes in expected remaining return
caused by an action) gives a naive range coefficient \(HR\), which
their concentration and stopping-time argument reduces to order \(R\)
up to logarithms. Their Proposition~3.1 also bounds expert-action
variance by the variance of the total return; a deterministic expert
has zero action variance even in a random environment.

Here \(\sigma_P(c)\) measures the variation of the specified trajectory
cost directly. For our binary examples it is zero, and the sharp
zero-cost bound already uses \(C_{\max}=T\) for the count and
\(C_{\max}=1\) for terminal error.
\citet{wang2025} survey variance-sensitive Hellinger bounds and the
role of loss functions in decision making.

The finite-class likelihood proof in \cref{app:finite-mle} makes the
optimization tolerance explicit.
Lecture~18 develops the imitation-learning results, including the expert
variance, parameter sharing, misspecification, and the sample-budget
meaning of horizon dependence.

The comparison-rule example in \cref{app:conditional-support} adapts
the rare-disagreement construction of \citet[Theorem~3.2]{gyorgy2025}
to conditional reverse KL. Its exact two-rule failure probability is
proved here.

\paragraph{Minimax optimality of the parameter-sharing rate.}
\phantomsection
\label{note:sharing-minimax}
The rate \(d(T)/n\) in \cref{prop:sharing-kl} is optimal up to
constants depending on \(a\). Write \(P_p\) for the length-\(T\)
grouped Bernoulli law with parameter vector \(p\). Then
\[
 \inf_{\widehat Q}\sup_{p\in[a,1-a]^{d(T)}}
 \E[\KL(P_p,\widehat Q)]\asymp_a\frac{d(T)}n,
\]
where the infimum ranges over all learners that use \(n\) independent
training sequences to output a probability law on length-\(T\) sequences.
The upper bound follows from \cref{prop:sharing-kl}.
For the lower bound, put independent two-point priors on the
parameters, with the two values near \(1/2\) separated by a small
multiple of \((nm_j)^{-1/2}\), depending only on \(a\).
At a test position in group \(j\),
the training data and preceding test tokens contain at most
\((n+1)m_j\le2nm_j\) observations of that parameter.
The two values remain hard to distinguish: Le Cam's testing bound
and Pinsker's inequality give expected conditional KL at least a
constant multiple of \(1/(nm_j)\). Summing over the \(m_j\) positions per group
and then over groups, using the KL chain rule, gives \(d(T)/n\).
\citet{yu1997} explains the two-point and coordinatewise testing
methods underlying this argument.

\paragraph{Loss normalization in supervised fine-tuning.}
\phantomsection
\label{note:sft-normalization}
\citet[Section~4.3, Eq.~(6)]{yao2025} distinguish averaging all predicted
token losses from averaging separately normalized response losses, and
compare these choices in supervised fine-tuning. \citet{hf2024loss}
explain how to implement token averaging across gradient-accumulation
batches: sum the losses and divide by the total number of prediction
targets included in the loss.

\paragraph{Statistical accuracy, optimization, and scaling laws.}
\phantomsection
\label{note:optimization-accuracy}
The strong-convexity certificate in \cref{sec:likelihood-pitfall} follows
by minimizing the quadratic lower bound in the definition of strong
convexity; \cref{ex:optimization-certificate} gives the calculation.
\citet[Section~9.1.2]{boyd2004} discuss this stopping criterion.
\citet[Theorem~4.7]{bottou2018} give the \(O(1/K)\) expected optimization
rate used in \cref{en:optimization-compute}, under smoothness, strong
convexity, controlled gradient noise, and suitable decreasing stepsizes.
Their Section~4.3 treats the weaker stationarity guarantees available
for general nonconvex objectives.

\citet[Sections~2--3]{bottou2007} study the
approximation--estimation--optimization tradeoff. Balancing optimization
tolerance against estimation error gives the statistical-accuracy
principle used in \cref{eq:statistical-accuracy}.
As data improve the statistical accuracy, finer optimization becomes
worthwhile.

The connection to language-model scaling laws is explicit in
\citet[Section~3.3 and Appendix~D.2]{hoffmann2022}. Renaming their
variables, they fit the empirical relation
\[
 L(m,B)\approx L_\infty+a\,m^{-\alpha}+b\,B^{-\beta},
\]
where \(m\) is parameter count, \(B\) is the number of training tokens,
and the constants are fitted across training runs.
Their last term combines finite-data effects and incomplete optimization
during a single pass. Increasing the training-token budget supplies
both observations and gradient updates.
Under a fixed compute budget, \citet{kaplan2020} find a regime in
which a larger model stopped before convergence outperforms smaller
models trained longer. This allocates computation between model size
and fitting accuracy.
These empirical laws describe average log loss. The cost comparisons
in \cref{sec:binary} ask what task-performance variation remains at
a given likelihood, and hence what additional information is useful
for choosing a generator.

\paragraph{Multi-token prediction and auxiliary heads.}
\phantomsection
\label{note:multistep-heads}
\citet[Section~2]{gloeckle2024} train several future-token heads on a
shared transformer representation and report improved downstream
performance in their experiments. These are separate prediction heads;
the architecture does not enforce equality with the autoregressive
model's multistep marginals. The exact loss in \cref{eq:multistep-loss}
and this practical auxiliary objective therefore need to be distinguished.
\citet[Section~2]{hinton2015} describe distillation by training a student
on a teacher's predictive probabilities. Using generator rollouts to
train a future-token head is a sampling-based version of this
prediction-matching idea.

\paragraph{Stability and concentration.}
\citet[Section~1]{gaubert2015} recall the Dobrushin ergodicity coefficient
and its characterization as the optimal TV contraction factor.
\citet{rudolf2018} develop perturbation bounds for Markov chains using
contraction. \Cref{thm:markov-stability} is the elementary finite-state
TV argument. The history-dependent comparison in
\cref{app:history-recovery} follows by replacing conditionals one
at a time and coupling the resulting continuations.
\citet[Section~3.1]{doukhan2007} use related summable Lipschitz
coefficients to study processes with infinite memory.
\citet[Sections~2.2.1 and~3.4.3]{raginsky2014} give
Hoeffding's lemma and the relative-entropy variational inequality.
Our exponential-tilting calculation applies these ideas to a specified
sequence cost. \Cref{prop:reverse-modulus} identifies its exact
reverse-KL conversion by constructing extremizing tilted laws.
The single-bit comparison in \cref{app:bound-comparison} uses the
sharp Bernoulli subgaussian proxy from \citet{schlemm2014}.

\paragraph{Moduli of continuity and approximation.}
\Cref{sec:discrepancy-comparison} applies the usual quantitative
continuity question to KL and expected cost. The fixed-KL endpoint
calculations give both conversion directions directly.
For the approximation-theory connection in \cref{app:moduli-analysis},
\citet[Theorem~1]{drucker2011} state Jackson's theorem in the form used
here and give proofs. Its modulus controls the uniform polynomial
approximation error as a function of degree.

\paragraph{Reverse-KL fitting.}
\label{note:reverse-fitting}
\citet[Section~2]{hinton2015} use cross-entropy with teacher probabilities,
equivalent to forward KL, for model compression through distillation.
\citet[Sections~1--2]{gu2024} motivate reverse-KL distillation by
prioritizing plausible student generations when the student cannot
represent the full teacher distribution. Their MiniLLM algorithm uses
policy gradients with practical stabilization; its reported benefits are
empirical. Their Sections~2.2--3.1 describe the generation scheme and
training data summarized in \hyperref[en:distillation]{the distillation endnote}.
\citet[Sections~2.2--2.3]{nowozin2016} develop adversarial
fitting through divergence lower bounds. Specializing their representation
to reverse KL and writing their critic as \(-e^g\) gives
\cref{eq:reverse-adversarial}.
\Cref{ex:reverse-variational} asks for a direct proof.
\Cref{app:reverse-optimization} derives the gradient for an evaluable
target and the cost-defined Gibbs target.

\begin{thebibliography}{99}
\bibitem[Bottou and Bousquet(2007)]{bottou2007}
L\'eon Bottou and Olivier Bousquet.
\newblock The Tradeoffs of Large Scale Learning.
\newblock \emph{Advances in Neural Information Processing Systems}, 20, 2007.
\newblock \url{https://leon.bottou.org/publications/pdf/nips-2007.pdf}.

\bibitem[Bottou et al.(2018)]{bottou2018}
L\'eon Bottou, Frank E. Curtis, and Jorge Nocedal.
\newblock Optimization Methods for Large-Scale Machine Learning.
\newblock \emph{SIAM Review}, 60(2):223--311, 2018.
\newblock \url{https://arxiv.org/abs/1606.04838}.

\bibitem[Boyd and Vandenberghe(2004)]{boyd2004}
Stephen Boyd and Lieven Vandenberghe.
\newblock \emph{Convex Optimization}.
\newblock Cambridge University Press, 2004. Section~9.1.2.
\newblock \url{https://web.stanford.edu/~boyd/cvxbook/}.

\bibitem[Doukhan and Wintenberger(2007)]{doukhan2007}
Paul Doukhan and Olivier Wintenberger.
\newblock Weakly dependent chains with infinite memory.
\newblock arXiv:0712.3231, 2007.
\newblock \url{https://arxiv.org/abs/0712.3231}.

\bibitem[Drucker and de Wolf(2011)]{drucker2011}
Andrew Drucker and Ronald de Wolf.
\newblock Uniform Approximation by (Quantum) Polynomials.
\newblock \emph{Quantum Information and Computation},
11(3\&4):215--225, 2011.
\newblock \url{https://arxiv.org/abs/1008.1599}.

\bibitem[Foster et al.(2024)]{foster2024}
Dylan J. Foster, Adam Block, and Dipendra Misra.
\newblock Is Behavior Cloning All You Need? Understanding Horizon
in Imitation Learning.
\newblock \emph{Advances in Neural Information Processing Systems}, 2024.
\newblock \url{https://arxiv.org/abs/2407.15007}.

\bibitem[Gaubert and Qu(2015)]{gaubert2015}
St\'ephane Gaubert and Zheng Qu.
\newblock Dobrushin's ergodicity coefficient for Markov operators on cones.
\newblock \emph{Integral Equations and Operator Theory}, 81(1):127--150, 2015.
\newblock \url{https://arxiv.org/abs/1307.4649}.

\bibitem[Gloeckle et al.(2024)]{gloeckle2024}
Fabian Gloeckle, Badr Youbi Idrissi, Baptiste Rozi\`ere,
David Lopez-Paz, and Gabriel Synnaeve.
\newblock Better \& Faster Large Language Models via Multi-token Prediction.
\newblock \emph{Proceedings of the 41st International Conference on Machine
Learning}, PMLR 235:15706--15734, 2024.
\newblock \url{https://proceedings.mlr.press/v235/gloeckle24a.html}.

\bibitem[Gu et al.(2024)]{gu2024}
Yuxian Gu, Li Dong, Furu Wei, and Minlie Huang.
\newblock MiniLLM: Knowledge Distillation of Large Language Models.
\newblock \emph{International Conference on Learning Representations}, 2024.
\newblock \url{https://proceedings.iclr.cc/paper_files/paper/2024/hash/8ac015d409635f196f9e3e9dcfb9a94e-Abstract-Conference.html}.

\bibitem[Gy\"orgy et al.(2025)]{gyorgy2025}
Andr\'as Gy\"orgy, Tor Lattimore, Nevena Lazi\'c, and Csaba Szepesv\'ari.
\newblock Beyond Statistical Learning: Exact Learning Is Essential
for General Intelligence.
\newblock arXiv:2506.23908, 2025.
\newblock \url{https://arxiv.org/abs/2506.23908}.

\bibitem[Hinton et al.(2015)]{hinton2015}
Geoffrey Hinton, Oriol Vinyals, and Jeff Dean.
\newblock Distilling the Knowledge in a Neural Network.
\newblock arXiv:1503.02531, 2015.
\newblock \url{https://arxiv.org/abs/1503.02531}.

\bibitem[Hoffmann et al.(2022)]{hoffmann2022}
Jordan Hoffmann et al.
\newblock Training Compute-Optimal Large Language Models.
\newblock arXiv:2203.15556, 2022.
\newblock \url{https://arxiv.org/abs/2203.15556}.

\bibitem[Kaplan et al.(2020)]{kaplan2020}
Jared Kaplan et al.
\newblock Scaling Laws for Neural Language Models.
\newblock arXiv:2001.08361, 2020.
\newblock \url{https://arxiv.org/abs/2001.08361}.

\bibitem[Nowozin et al.(2016)]{nowozin2016}
Sebastian Nowozin, Botond Cseke, and Ryota Tomioka.
\newblock f-GAN: Training Generative Neural Samplers using Variational
Divergence Minimization.
\newblock \emph{Advances in Neural Information Processing Systems}, 29, 2016.
\newblock \url{https://arxiv.org/abs/1606.00709}.

\bibitem[Raginsky and Sason(2014)]{raginsky2014}
Maxim Raginsky and Igal Sason.
\newblock \emph{Concentration of Measure Inequalities in Information
Theory, Communications and Coding}.
\newblock Second edition, 2014.
\newblock \url{https://arxiv.org/abs/1212.4663}.

\bibitem[Rudolf and Schweizer(2018)]{rudolf2018}
Daniel Rudolf and Nikolaus Schweizer.
\newblock Perturbation theory for Markov chains via Wasserstein distance.
\newblock \emph{Bernoulli}, 24(4A):2610--2639, 2018.
\newblock \url{https://doi.org/10.3150/17-BEJ938}.

\bibitem[Schlemm(2014)]{schlemm2014}
Eckhard Schlemm.
\newblock The Kearns--Saul inequality for Bernoulli and Poisson-binomial
distributions.
\newblock arXiv:1405.4496, 2014.
\newblock \url{https://arxiv.org/abs/1405.4496}.

\bibitem[Wang et al.(2025)]{wang2025}
Kaiwen Wang, Nathan Kallus, and Wen Sun.
\newblock The Central Role of the Loss Function in Reinforcement Learning.
\newblock arXiv:2409.12799, 2025.
\newblock \url{https://arxiv.org/abs/2409.12799}.

\bibitem[Yu(1997)]{yu1997}
Bin Yu.
\newblock Assouad, Fano, and Le Cam.
\newblock In David Pollard, Erik Torgersen, and Grace L. Yang, editors,
\emph{Festschrift for Lucien Le Cam: Research Papers in Probability
and Statistics}, pages 423--435. Springer, 1997.
\newblock \url{https://doi.org/10.1007/978-1-4612-1880-7_29}.

\bibitem[Hugging Face(2024)]{hf2024loss}
Hugging Face.
\newblock Fixing gradient accumulation.
\newblock October 16, 2024.
\newblock \url{https://huggingface.co/blog/gradient_accumulation}.

\bibitem[Yao et al.(2025)]{yao2025}
Yongqiang Yao, Jingru Tan, Kaihuan Liang, Feizhao Zhang, Jiahao Hu,
Shuo Wu, Yazhe Niu, Ruihao Gong, Dahua Lin, and Ningyi Xu.
\newblock Hierarchical balance packing: Towards efficient supervised
fine-tuning for long-context LLM.
\newblock NeurIPS, 2025. arXiv:2503.07680v3.
\newblock \url{https://arxiv.org/abs/2503.07680v3}.

\end{thebibliography}

\clearpage

\section*{Glossary}
\phantomsection
\label{sec:glossary}

\begin{description}[style=nextline,leftmargin=1.5em,labelindent=0pt,
  itemsep=0.15em,parsep=0pt]
\item[Sequence model; autoregressive factorization]
A joint law on token sequences, expressed as a product of next-token
conditionals. Every finite sequence law has such a factorization.
\item[Maximum likelihood; teacher forcing]
Maximum likelihood minimizes empirical negative log probability.
Teacher forcing evaluates each conditional using the observed prefix.
Generation instead feeds sampled tokens back into the predictor.
\item[Positive-data fine-tuning]
Continuing model fitting on examples selected to demonstrate desired
behavior. The idealization here uses sequences of cost zero.
\item[Cost and reward]
A numerical assessment of a generated response or sequence.
The task objective is its expected cost, or expected reward.
\item[Cost discrepancy; moduli of continuity]
The pseudometric \(d_c(P,Q)=|\E_P[c]-\E_Q[c]|\) compares
one expected cost. The modulus \(\omega_{D\to c}\) asks how much cost
discrepancy a KL budget permits. The converse \(\omega_{c\to D}\) asks
how large KL can be at a given cost-discrepancy budget. The upper and
lower endpoints of the fixed-KL cost ranges determine these two moduli.
Taking the supremum over reference laws gives a uniform conversion.
\item[Hellinger distance (appendix)]
A distributional distance used in the proofs of the variance-sensitive
cost and finite-sample likelihood bounds; defined in \cref{app:hellinger}.
\item[Optimization error and tolerance]
The error \(E(\widehat Q)\) is the model's excess empirical per-token
loss above the infimum in the class. A tolerance \(\eta\) bounds this
error and permits per-sequence
excess \(T\eta\). The average can place this excess unevenly across positions.
\item[Total variation]
\(\tv{P-Q}\) is the largest discrepancy in event probabilities.
For costs in \([0,C_{\max}]\), it bounds expectation discrepancy
multiplied by \(C_{\max}\).
\item[Forward and reverse KL]
\(\KL(P, Q)\) averages log density ratios under the data law;
\(\KL(Q, P)\) averages them under the model law.
The sequence chain rules inherit these different prefix distributions.
\item[Markov process; transition kernel]
A process whose current state determines the conditional next-state
law. \(K(x,y)\) gives its transition probabilities.
\item[Marginal law]
The distribution of a selected part of a random sequence.
The marginal at position \(t\) assigns probability \(P(X_t=x)\)
to \(x\), summing over the other positions.
\item[Dobrushin ergodicity coefficient; contraction margin]
The coefficient \(\rho(K)\) is the largest TV distance between two rows
of \(K\), and its optimal TV contraction factor.
If a kernel contracts TV by at most \(\rho<1\), its margin is
\(1-\rho\). The local-error bound is amplified by its reciprocal.
\item[Recovery]
In the binary example, the probability of returning from an erroneous
state to zero. It controls the duration of an erroneous run.
\item[Subgaussian cost]
A cost whose centered log exponential moment is at most
\(\lambda^2v^2/2\). The quantity \(v^2\) is a variance proxy;
\(v_P(c)\) is the smallest valid proxy under \(P\).
It satisfies \(\sigma_P(c)^2\le v_P(c)\le C_{\max}^2/4\)
for \(0\le c\le C_{\max}\).
\item[\(k\)-step prediction loss]
\(\cL_k(Q)\) scores a token \(k\) positions ahead from the observed
prefix, summing over the model's intervening tokens.
\item[Critic; adversarial fitting]
A critic scores sequences to maximize a divergence lower bound.
Adversarial fitting alternates critic maximization with generator
minimization; a restricted critic may leave a gap to the full divergence.
\end{description}


\clearpage

\section{Exercises}
\label{sec:exercises}

These exercises follow the lecture from KL and task cost to sequence
length, fitting, recovery, and persistent memory.
\exerciseinstructions
\setcounter{courseexercise}{0}

% Shared exercise chunk: l4-kl-chain
\begin{exercisechunk}
\item \textbf{The KL chain rule.}
\label[exercise]{ex:kl-chain}
Let \(P,Q\) be laws on \([N]^T\), with \(Q(x)>0\) whenever \(P(x)>0\).
\begin{enumerate}[(a),beginpenalty=10000]
\item Use the probability chain rule to write
\(\log(P(X_{1:T})/Q(X_{1:T}))\) as a sum.
\item Take expectation under \(P\), and then condition on \(X_{<t}\)
in each term, to prove the KL chain rule in \cref{eq:population-kl}.
\item State the reversed identity, with its own support condition.
Which prefix distribution appears in each direction?
\item Let \(P,Q\) now be laws on \(\mathcal U\times[N]^T\), where
\(\mathcal U\) is a finite prompt set, with the same prompt marginal
\(\mu\). Prove
\[
 \KL(P,Q)=\sum_{u:\,\mu(u)>0}\mu(u)
 \KL(P(\,\cdot\mid u),Q(\,\cdot\mid u)).
\]
\end{enumerate}
\end{exercisechunk}

% Shared exercise chunk: l4-positive-data
\begin{exercisechunk}
\item \textbf{Zero-cost data and forward KL.}
\label[exercise]{ex:positive-data}
Let \(P,Q\) be laws on a finite set \(\cX\), and let
\(c:\cX\to[0,C_{\max}]\), with \(C_{\max}>0\) and \(\E_P[c]=0\).
\begin{enumerate}[(a),beginpenalty=10000]
\item Prove the bound in \cref{eq:positive-data-bound}.
\item Let \(\kappa\ge0\) and let \(R\) be supported on
\(\{x:c(x)=C_{\max}\}\). Show that
\(Q=e^{-\kappa}P+(1-e^{-\kappa})R\) attains the bound.
\item Can a smaller bound hold for all such laws and costs using only
\(C_{\max}\) and \(\KL(P,Q)\)? Explain using part (b).
\end{enumerate}
\begin{hints}
For (a), take \(A=\{x:c(x)=0\}\) and compare \(P(A)\) with \(Q(A)\).
\end{hints}
\end{exercisechunk}

% Shared exercise chunk: l4-reverse
\begin{exercisechunk}
\item \textbf{Concentration and reverse KL.}
\label[exercise]{ex:reverse}
Use the notation of \cref{sec:cost-concentration}.
\begin{enumerate}[(a),beginpenalty=10000]
\item Prove \cref{prop:entropy-cost} by exponential tilting, and
then deduce \cref{thm:reverse-cost}, including the case \(v=0\).
\item Let \(c=T^{-1}\sum_{t=1}^T f_t(X_t)\), with independent
\([0,1]\)-valued summands under \(P\). Prove
\(v_P(c)\le1/(4T)\), and deduce, for every law \(Q\),
\begin{equation}
 |\E_Q[c]-\E_P[c]|\le
 \sqrt{\frac{\KL(Q, P)}{2T}}.
 \label{eq:reverse-average}
\end{equation}
What growth of reverse KL still permits vanishing discrepancy?
\item Let \(P_\delta=\operatorname{Bernoulli}(\delta)^{\otimes T}\),
\(0<\delta<1\), and set \(m=\E_Q[c_{\mathrm{cnt}}]/T\) for the error-count
cost. Prove
\[
 \KL(Q, P_\delta)
 =T\,\operatorname{kl}(m,\delta)+\KL(Q, P_m),
\]
where \(P_m=\operatorname{Bernoulli}(m)^{\otimes T}\) and
\(\operatorname{kl}(a,b)=a\log(a/b)+(1-a)\log((1-a)/(1-b))\).
Treat \(m=0,1\) as well.
\item Which generator has the smallest reverse KL among laws with a
specified expected count?
\end{enumerate}

\begin{hints}
For (a), normalize \(P(x)e^{\lambda(c(x)-\E_P[c])}\) to obtain a
probability law, then use nonnegativity of KL. For (b), apply
Hoeffding's lemma to each independent summand.
\end{hints}
\end{exercisechunk}

% Shared exercise chunk: l4-discrepancy-comparison
\begin{exercisechunk}
\item \textbf{Comparing fitting and task discrepancies.\suppmark}
\label[exercise]{ex:discrepancy-comparison}
Use the definitions in \cref{sec:discrepancy-comparison}.
For parts (b)--(e), let \(P=\delta_{0^T}\), \(T\ge2\), with
\(c_{\mathrm{cnt}}(x)=\sum_{t=1}^T x_t\) and
\(c_{\mathrm{fin}}(x)=x_T\).
\begin{enumerate}[(a),beginpenalty=10000]
\item Prove that \(d_c\) is a pseudometric. Give distinct laws with
\(d_c(P,Q)=0\), first for the binary count and then for another cost.
\item Use the fixed-KL intervals
in \cref{eq:binary-kl-level,eq:terminal-kl-level} to prove
\begin{align}
 \omega_{D\to c_{\mathrm{cnt}}}(\kappa;P)
 &=T(1-e^{-\kappa}),\label{eq:binary-moduli}\\
 \omega_{D\to c_{\mathrm{fin}}}(\kappa;P)
 &=1-e^{-\kappa},\qquad \kappa\ge0.
 \label{eq:terminal-moduli}
\end{align}
Which of $Q_-$ and $Q_+$ attains these bounds?
\item Calculate the reverse conversions:
\begin{align}
 \omega_{c_{\mathrm{cnt}}\to D}(s;P)
 &=\begin{cases}-\log(1-s),&0\le s<1,\\
                 +\infty,&s\ge1,\end{cases}
 \label{eq:binary-converse-modulus}\\
 \omega_{c_{\mathrm{fin}}\to D}(s;P)
 &=+\infty,\qquad s\ge0.
 \label{eq:terminal-converse-modulus}
\end{align}
Exhibit laws attaining each finite bound and laws with infinite KL
where appropriate. Explain why small expected count forces small
KL, whereas perfect final answers allow arbitrarily large KL.
\item Fix an expected count \(s\in[0,1)\). Prove the sharp bounds
\[
 -\log(1-s/T)\le\KL(P, Q)\le-\log(1-s).
\]
Give a law attaining each endpoint and relate them to part (c).
\item Replace the count by the average count \(c_{\mathrm{cnt}}/T\).
Recompute both moduli. How does normalization change their scales
and the ratio between the largest and smallest costs at a fixed
positive KL?
\item Verify that the cost in \cref{prop:forward-variance} may have
zero variance under \(P\) and positive expectation under \(Q\).
Which term in the bound accounts for this?
\end{enumerate}
\end{exercisechunk}

% Shared exercise chunk: l4-binary-count
\begin{exercisechunk}
\item \textbf{Repeated errors and their accumulation.}
\label[exercise]{ex:binary-count}
Compare isolated errors with errors that persist. Fix
\(T\ge2\), \(P=\delta_{0^T}\), \(0<\varepsilon<1\), and \(X_0=0\).
Use \(c_{\mathrm{cnt}}(x)=\sum_{t=1}^T x_t\) and
\(c_{\mathrm{fin}}(x)=x_T\).
The Markov generators \(Q_{\mathrm{ind}},Q_{\mathrm{bad}}\) use
\[
 \begin{array}{c|cc}
 &Q(X_{t+1}=1\mid X_t=0)&Q(X_{t+1}=1\mid X_t=1)\\ \hline
 Q_{\mathrm{ind}}&\varepsilon&\varepsilon\\
 Q_{\mathrm{bad}}&\varepsilon&1
 \end{array}
 \qquad(0\le t<T).
\]
For \(P\), use the kernel that returns zero from either state.
\begin{enumerate}[(a),beginpenalty=10000]
\item Show that both generators have conditional KL
\(-\log(1-\varepsilon)\) at every data prefix, and hence
\begin{equation}
 \KL(P, Q_{\mathrm{ind}})
 =\KL(P, Q_{\mathrm{bad}})
 =T\log\frac1{1-\varepsilon}.
 \label{eq:binary-kl}
\end{equation}
Compare this constant positionwise loss with the loss profile of
$Q_-$ and $Q_+$. Compute the probability that the first one occurs at
position \(s\), and the probability of no ones.
\item Prove
\begin{align}
 \E_{Q_{\mathrm{ind}}}[c_{\mathrm{cnt}}]&=T\varepsilon,
 \label{eq:ind-cost}\\
 \E_{Q_{\mathrm{bad}}}[c_{\mathrm{cnt}}]
 &=\sum_{t=1}^T[1-(1-\varepsilon)^t]
 =T-\frac{(1-\varepsilon)[1-(1-\varepsilon)^T]}{\varepsilon}.
 \label{eq:bad-cost}
\end{align}
Obtain the second identity both from marginal error probabilities
and by conditioning on the first error position. Also prove
\begin{equation}
 \E_{Q_{\mathrm{ind}}}[c_{\mathrm{fin}}]=\varepsilon,\qquad
 \E_{Q_{\mathrm{bad}}}[c_{\mathrm{fin}}]=1-(1-\varepsilon)^T.
 \label{eq:terminal-cost}
\end{equation}
\item Prove the finite bounds
\[
 \frac{\varepsilon T(T+1)}2
 -\frac{\varepsilon^2T(T+1)(T-1)}6
 \le\E_{Q_{\mathrm{bad}}}[c_{\mathrm{cnt}}]
 \le\frac{\varepsilon T(T+1)}2.
\]
Deduce \(\E_{Q_{\mathrm{bad}}}[c_{\mathrm{cnt}}]
\sim\varepsilon T(T+1)/2\) when \(T\varepsilon\to0\).
Explain why the expected count can grow quadratically with \(T\)
in this regime.
\item Show that both sequence TVs equal \(1-(1-\varepsilon)^T\).
For either generator, verify the bounds
\begin{align}
 \E_Q[c_{\mathrm{cnt}}]&\le T[1-(1-\varepsilon)^T],
 &&\text{using exact TV},\label{eq:binary-tv-bound}\\
 \E_Q[c_{\mathrm{cnt}}]&\le
 T\min\left\{1,\sqrt{\frac T2\log\frac1{1-\varepsilon}}\right\},
 &&\text{using Pinsker}.\label{eq:binary-pinsker-bound}
\end{align}
Compare these with the exact counts at
\(T=1000,\varepsilon=10^{-3}\), and their leading terms when
\(T\varepsilon\to0\). Which information about error duration
does the TV-to-count bound discard?
\item Set \(\varepsilon=1-e^{-\kappa/T}\) for fixed \(\kappa>0\).
Compute the two terminal costs and the ratio
\(\E_{Q_{\mathrm{bad}}}[c_{\mathrm{fin}}]/
\E_{Q_{\mathrm{ind}}}[c_{\mathrm{fin}}]\). Explain why this ratio is of order \(T\) even though the terminal cost range is
fixed. Which model attains the terminal-cost upper endpoint?
\end{enumerate}

\end{exercisechunk}

% Shared exercise chunk: l4-positionwise-kl
\begin{exercisechunk}
\item \textbf{Positionwise and average KL bounds.}
\label[exercise]{ex:positionwise-kl}
Let \(P=\delta_{0^T}\), \(T\ge2\), and use the counting and terminal
costs from \cref{ex:binary-count}. Compare a bound on each
next-token KL with a bound on their average. Fix \(\eta\ge0\).
\begin{enumerate}[(a),beginpenalty=10000]
\item Suppose a generator \(Q\) satisfies
\[
 \E\!\left[\KL\bigl(P(\,\cdot\mid X_{<t})
       , Q(\,\cdot\mid X_{<t})\bigr)\right]\le\eta
 \quad(t=1,\ldots,T),\qquad X_{1:T}\sim P.
\]
Show that \(Q(X_{1:t}=0^t)\ge e^{-t\eta}\), and deduce
\begin{equation}
 \E_Q[c_{\mathrm{cnt}}]\le\sum_{t=1}^T(1-e^{-t\eta}),
 \qquad
 \E_Q[c_{\mathrm{fin}}]\le1-e^{-T\eta}.
 \label{eq:positionwise-count}
\end{equation}
Show that \(Q_{\mathrm{bad}}\) attains both bounds when
\(\varepsilon=1-e^{-\eta}\).
\item Suppose instead that only \(\KL(P, Q)/T\le\eta\).
Use \cref{eq:binary-kl-level,eq:terminal-kl-level} to obtain the
corresponding bounds. Compare their small-\(T\eta\) leading terms
with \cref{eq:positionwise-count}. What per-token accuracy suffices
to guarantee terminal error at most \(s\in(0,1)\)?
\end{enumerate}

\end{exercisechunk}

% Shared exercise chunk: l4-sharing
\begin{exercisechunk}
\item \textbf{Estimating shared parameters.\suppmark}
\label[exercise]{ex:sharing}
Use the independent-bit models and pooled estimator from
\cref{sec:parameter-sharing}. All expectations below average over the
training observations. Fix \(0<a<1/2\).
\begin{enumerate}[(a),beginpenalty=10000]
\item Let \(B\sim\operatorname{Binomial}(L,p)\), with \(L\ge1\).
For \(0<p<1\), show that the ordinary empirical frequency
\(\bar p=B/L\) satisfies \(\E[\kl(p,\bar p)]=\infty\).
Explain why adding one count of each bit gives
\(\widehat p=(B+1)/(L+2)\), and why this assigns positive probability
to both bits. To obtain a finite-sample bound for this choice, prove
that, for \(0<p<1\),
\[
 \E\!\left[\frac1{B+1}\right]
 =\frac{1-(1-p)^{L+1}}{(L+1)p}.
\]
Apply the same identity to \(L-B\), then use Jensen's inequality
for the logarithm to show
\[
 \E[\kl(p,\widehat p)]
 \le\log\frac{L+2}{L+1}
 \le\frac1{L+1}.
\]
Verify the bound directly for \(p=0,1\).
\item Apply part (a) with \(L=nm_j\) to bound \(\kappa_n(T)\) from above.
For the lower bound, calculate
\[
 \E[(\widehat p-p)^2]
 =\frac{Lp(1-p)+(1-2p)^2}{(L+2)^2}.
\]
Use Pinsker's inequality and \(L/(L+2)\ge1/3\) to show that
\[
 \frac{2a(1-a)}9\frac{d(T)}n
 \le\kappa_n(T)\le\frac{d(T)}n
 \qquad\text{when }p_j\in[a,1-a].
\]
Explain why group sizes cancel from both bounds.
\item For fixed \(p\in(0,1)\), prove
\(L\,\E[\kl(p,\widehat p)]\to1/2\) as \(L\to\infty\).
Deduce \cref{eq:shared-bernoulli,eq:unshared-bernoulli}, keeping track
of which of \(n,T\) is fixed in each limit.
Make the expansion uniform for \(p\in[a,1-a]\) and prove
\cref{eq:sharing-leading} when \(n\min_jm_j\to\infty\).
\end{enumerate}
\begin{hints}
For (a), use \(1/(B+1)=\int_0^1u^B\,du\).
For (c), expand \(\kl(p,q)\) near \(q=p\) and use a binomial tail
bound for the remaining event.
\end{hints}
\end{exercisechunk}

% Shared exercise chunk: l4-sharing-horizon
\begin{exercisechunk}
\item \textbf{Sequence length and the training budget.\suppmark}
\label[exercise]{ex:sharing-horizon}
Use the grouped Bernoulli model and smoothed estimator from
\cref{sec:parameter-sharing}. Assume \(p_j\in[a,1-a]\), with
\(0<a<1/2\).
\begin{enumerate}[(a),beginpenalty=10000]
\item Keep \(n\) fixed and suppose \(p_t\in[a,1-a]\).
Use \cref{ex:sharing}(b) to give explicit upper and lower bounds proportional to
\(T/n\) for the unshared model's expected sequence KL.
Compare this growth with the shared model's limit in
\cref{eq:shared-bernoulli}.
\item Use one new parameter for each successive block of lengths
\(1,2,3,\ldots\). Compute \(d(T)\) at complete-block endpoints and
interpret \cref{eq:sharing-rate,eq:sharing-leading}. How can normalized
KL improve while sequence KL grows?
\item Fix the total training-token budget \(L=nT\) and use
\(\kappa'_L(T)=\kappa_{L/T}(T)\) to repeat the shared and unshared
leading-order comparisons. Explain why increasing \(T\) worsens
expected sequence KL in the shared model. Compare the numbers of
observations per parameter in the two models. Translate the expected-KL bounds into
bounds on the discrepancy of a cost in \([0,1]\).
\end{enumerate}

\end{exercisechunk}

% Shared exercise chunk: l4-sharing-costs
\begin{exercisechunk}
\item \textbf{Task costs at a given learning rate.\suppmark}
\label[exercise]{ex:sharing-costs}
Prove \cref{prop:sharing-full-cost} using the following construction.
At each horizon \(T\), partition
the first \(T\) positions into \(\lceil d(T)\rceil\) nonempty groups.
Let the \(U_t\) be independent fair bits and \(V_t=0\) under \(P\).
Use terminal cost \(V_T\) and counting cost \(\sum_{t=1}^T V_t\).
Write \(P_T\) and \(P_T^U\) for the laws of \((U,V)_{1:T}\) and
\(U_{1:T}\), respectively.
If \(\widehat P^U_T\) is the pooled
smoothed estimate of the law of \(U_{1:T}\), set
\[
 r_T=\E\!\left[\KL(P^U_T,\widehat P^U_T)\right]
\]
and, for any law $R$ supported on $v_{1:T}\ne0^T$, define
\[
 \widehat Q_R
 =e^{-r_T}(\widehat P^U_T\otimes\delta_{0^T})
  +(1-e^{-r_T})R.
\]
\begin{enumerate}[(a),beginpenalty=10000]
\item Show that \(r_T\) is determined by \(n,T\) and the chosen
partition, so the mixture weight can be used by the estimator.
\item Compute \(\E[\KL(P_T,\widehat Q_R)]\) and show that it has
order \(d(T)/n\), independently of \(R\).
\item Determine the full attainable intervals for the expected terminal
and counting costs as \(R\) varies. Identify laws attaining the endpoints.
\end{enumerate}

\end{exercisechunk}

% Shared exercise chunk: l4-bernoulli-identification
\begin{exercisechunk}
\item \textbf{Identifying one of two Bernoulli models.\suppmark}
\label[exercise]{ex:bernoulli-identification}
The data law is one of two infinite i.i.d.\ Bernoulli processes, with
parameter \(p\in\{1/4,3/4\}\). Observe \(n\) independent length-\(T\)
sequences. Choose \(\widehat p=3/4\) if the empirical mean of their
\(nT\) bits is at least \(1/2\), and choose \(1/4\) otherwise.
Let \(\widehat Q\) be the i.i.d.\ model with parameter \(\widehat p\).
\begin{enumerate}[(a),beginpenalty=10000]
\item Prove that the probability of choosing the wrong parameter is
at most \((\sqrt3/2)^{nT}\le e^{-nT/8}\).
\item Compute the sequence KL when the selected parameter is wrong,
and deduce
\[
 \kappa_n(T)\le\frac{T\log3}{2}e^{-nT/8}\longrightarrow0
 \qquad(T\to\infty,\ n\text{ fixed}).
\]
\item Compare this with estimating an arbitrary \(p\in(0,1)\) in
\cref{eq:shared-bernoulli}. What allows the two-candidate learner
to identify the entire model exactly?
\end{enumerate}

\begin{hints}
For (a), let \(B\) be the number of observed ones and apply Markov's
inequality to \(3^{B-nT/2}\) when \(p=1/4\).
\end{hints}
\end{exercisechunk}

% Shared exercise chunk: l4-copying-context
\begin{exercisechunk}
\item \textbf{Copying with a bounded context.\suppmark}
\label[exercise]{ex:copying-context}
Let \(P\) generate \(B_1B_1B_2B_2\cdots\), where the blocks
\(B_m\) are independent and uniform on \(\{0,1\}^m\).
For an integer \(w\ge0\), let \(\mathcal Q_w\) be all
infinite-sequence models whose next-token distribution can depend
on the position and the last \(w\) tokens only.
\begin{enumerate}[(a),beginpenalty=10000]
\item In a copied block of length \(m>w\), show that the required
bit is fair given the available context, although the full prefix
determines it. Deduce expected conditional KL at least \(\log2\)
at each such position.
\item At \(T=M(M+1)\), for \(M>w\), prove
\begin{equation}
 \inf_{Q\in\mathcal Q_w}\KL(P_{1:T}, Q_{1:T})
 =\frac{\log2}{2}\bigl(T-w(w+1)\bigr).
 \label{eq:copying-context}
\end{equation}
Specify a predictor attaining equality.
\item Explain why more training data cannot remove this error for
a fixed \(w\). How must \(w\) grow to predict every copied bit exactly?
\item Show that an exact recurrent generator must distinguish
\(2^m\) states after generating a fresh length-\(m\) block, and hence
needs at least \(m\) bits of state information.
Compare the description length of the copying rule with its memory
requirement. Relate the distinction to recalling earlier definitions
or intermediate results in a mathematical argument.
\end{enumerate}
\end{exercisechunk}

% Shared exercise chunk: l4-positionwise-loss
\begin{exercisechunk}
\item \textbf{A missed dependency at the final prediction.\suppmark}
\label[exercise]{ex:positionwise-loss}
Let \(T\ge2\). Under \(P_{\mathrm{par}}\), the first \(T-1\) bits
are independent fair coins and the last is their parity:
\(X_T=(\sum_{t<T}X_t)\bmod2\). Under \(Q_{\mathrm{coin}}\), all
\(T\) bits are independent fair coins. The task cost is
\[
 c_{\mathrm{par}}(x_{1:T})
 =\I\!\left\{x_T\ne\Bigl(\sum_{t<T}x_t\Bigr)\bmod2\right\}.
\]
\begin{enumerate}[(a),beginpenalty=10000]
\item Compute the next-token KL at every position and prove
\begin{equation}
 \begin{gathered}
 \KL(P_{\mathrm{par}}, Q_{\mathrm{coin}})=\log2,\qquad
 \frac{\KL(P_{\mathrm{par}}, Q_{\mathrm{coin}})}T
 =\frac{\log2}{T},\\
 \E_{P_{\mathrm{par}}}[c_{\mathrm{par}}]=0,\qquad
 \E_{Q_{\mathrm{coin}}}[c_{\mathrm{par}}]=\frac12.
 \end{gathered}
 \label{eq:parity-final-error}
\end{equation}
Verify that \cref{eq:positive-data-bound} is attained.
Compare the effect of normalization with \cref{eq:endpoint-normalization}.
\item For the model class \(\{P_{\mathrm{par}},Q_{\mathrm{coin}}\}\)
and any sample from \(P_{\mathrm{par}}\), prove that returning
\(Q_{\mathrm{coin}}\) has empirical per-token optimization gap
\((\log2)/T\). What has been left unlearned when this gap is accepted?
\item Now use the class of all independent-bit laws
\(Q=\bigotimes_{t=1}^T\operatorname{Bernoulli}(q_t)\).
Show that every bit is marginally fair under \(P_{\mathrm{par}}\),
and deduce that \(Q_{\mathrm{coin}}\) uniquely minimizes population
log loss in this class. Explain why the same prediction failure is
now approximation error.
\item If \(\KL(P, Q)/T\le\eta\), what is the largest possible
final-position contribution
\[
 \E\!\left[\KL\bigl(P(\,\cdot\mid X_{<T})
       , Q(\,\cdot\mid X_{<T})\bigr)\right],
 \qquad X_{1:T}\sim P,
\]
in \cref{eq:population-kl}? Give an example attaining it.
Compare the information supplied by the average and by the individual
positionwise losses.
\item Define the running parity \(Y_t=(\sum_{j=1}^t X_j)\bmod2\).
Describe its Markov kernels under \(P_{\mathrm{par}}\) and
\(Q_{\mathrm{coin}}\). Apply \cref{thm:markov-stability} to this
state, for which \(c_{\mathrm{par}}=Y_T\). Show that the fitted
kernels have contraction factor zero, all earlier local discrepancies
vanish, and the final local discrepancy is \(1/2\).
Why does contraction leave this final error unchanged?
\end{enumerate}
\end{exercisechunk}

% Shared exercise chunk: l4-optimization-certificate
\begin{exercisechunk}
\item \textbf{Certifying an optimization tolerance.\suppmark}
\label[exercise]{ex:optimization-certificate}
Let \(f:\R^d\to\R\) be differentiable, with
\(f_\star=\inf_{v\in\R^d}f(v)>-\infty\).
For \(f(\theta)=\widehat L(Q_\theta)\) and
\(\mathcal Q=\{Q_\theta:\theta\in\R^d\}\), the gap
\(f(\theta)-f_\star\) is the optimization error in
\cref{eq:approximate-fit-error}.
\begin{enumerate}[(a),beginpenalty=10000]
\item Suppose \(f\) is \emph{\(\mu\)-strongly convex}, meaning that
for a known \(\mu>0\) and every \(\theta,v\in\R^d\),
\[
 f(v)\ge f(\theta)+\nabla f(\theta)^\top(v-\theta)
                 +\frac\mu2\|v-\theta\|_2^2.
\]
Prove
\[
 f(\theta)-f_\star\le\frac{\|\nabla f(\theta)\|_2^2}{2\mu}.
\]
For a desired per-sequence tolerance \(s>0\), give a sufficient
gradient-norm threshold that certifies
\(T\bigl(f(\theta)-f_\star\bigr)\le s\).
\item Take \(d=1\) and \(f(\theta)=(\theta^2-1)^2\).
Compute \(f'(0)\) and \(f(0)-f_\star\). Explain why a small gradient
alone does not certify a small optimization error for a general loss.
\end{enumerate}
\begin{hints}
For (a), minimize the quadratic lower bound over \(v-\theta\).
\end{hints}
\end{exercisechunk}

% Shared exercise chunk: l4-fitting-support
\begin{exercisechunk}
\item \textbf{Likelihood accuracy, data, and support.\suppmark}
\label[exercise]{ex:fitting-support}
Parts (a)--(c) use the repeated-error generators from
\cref{ex:binary-count} and all-zero training sequences.
\begin{enumerate}[(a),beginpenalty=10000]
\item Compute both empirical per-token losses. Include
\(P=\delta_{0^T}\) in the model class and show that
\(\varepsilon=1-e^{-\eta}\) makes both generators acceptable at
optimization tolerance \(\eta>0\).
\item Compute their terminal and counting costs and their leading
terms as \(T\eta\to0\). Give a sufficient order of \(\eta\) in
\(T\) to keep each expected count bounded.
\item Remove \(P\) and use a finite collection of paired generators
at distinct positive values of \(\varepsilon\). Identify all
empirical minimizers. Explain why exact fitting leaves a tie between
models of different task cost.
\item Specialize \cref{thm:finite-mle} to zero-cost data, first for
\(c\in[0,T]\) and then for \(c\in[0,1]\). Express each guarantee
using the number of sequences \(n\), then using the token budget
\(L=nT\).
\item In \cref{prop:conditional-support}, choose a prompt on which
the two candidate rules disagree. Show that its expert label identifies
the correct rule. Compare the number of such queries with the number
of random demonstrations needed for failure probability at most
\(\delta\in(0,1/2)\).
\end{enumerate}
\end{exercisechunk}

% Shared exercise chunk: l4-multistep-benefit
\begin{exercisechunk}
\item \textbf{Two-step prediction and recovery.}
\label[exercise]{ex:multistep-benefit}
Use the losses in \cref{eq:multistep-loss}.
\begin{enumerate}[(a),beginpenalty=10000]
\item For \(Q_-,Q_+\) from \cref{sec:endpoint-generators}, with
\(T\ge2\) and \(\kappa>0\), compute \(\cL_1\) and \(\cL_2\)
on \(0^T\).
\item Which model minimizes \(\cL_1+\lambda\cL_2\) over this pair
when \(\lambda>0\), and by how much is its loss smaller?
Compare the models' expected terminal costs \(\E_Q[c_{\mathrm{fin}}]\).
\end{enumerate}
\end{exercisechunk}

% Shared exercise chunk: l4-multistep
\begin{exercisechunk}
\item \textbf{Computing further-ahead predictions.}
\label[exercise]{ex:multistep}
Use the autoregressive model in \cref{eq:multistep-marginal}. Fix
\(x_{\le t}\), \(k\ge2\) with \(t+k\le T\), and \(a\in[N]\).
\begin{enumerate}[(a),beginpenalty=10000]
\item Derive \cref{eq:multistep-marginal}.
Count its terms and the number of conditional evaluations if each term
is evaluated separately. Explain why sharing prefixes still leaves
exponentially many distinct prefixes in a general model.
\item Which conditional probability is trained if the observed
intervening tokens are supplied as context instead of being marginalized?
\item Suppose \(t\ge1\) and the model is first-order Markov on \([N]\)
with known transition matrix \(K\). Give a recursion that computes the
entire \(k\)-step distribution in \(O(kN^2)\) arithmetic operations.
Why can this recursion merge histories?
\item A separate head \(h_{\theta,k}(\,\cdot\mid x_{\le t})\) can be
trained from pairs \((X_{\le t},X_{t+k})\). State the additional
equality needed for its probabilities to be the original sequence
model's \(k\)-step marginals. Does perfect prediction on the observed
prefix--target pairs imply this equality? Give an example or a proof.
\end{enumerate}
\end{exercisechunk}

% Shared exercise chunk: l4-kl-directions
\begin{exercisechunk}
\item \textbf{Which errors do the KL directions penalize?}
\label[exercise]{ex:kl-directions}
\label[exercise]{prop:kl-directions}
Let \(P,R\) be laws on a finite domain \(\cX\).
\begin{enumerate}[(a),beginpenalty=10000]
\item For \(0<\varepsilon<1\), set \(Q=(1-\varepsilon)P+\varepsilon R\).
For \(X\sim P\), \(Z\sim R\), and \(Y\sim Q\), prove
\begin{equation}
 \KL(P,Q)\le-\log(1-\varepsilon),\qquad
 \E[c(Y)]=(1-\varepsilon)\E[c(X)]+\varepsilon\E[c(Z)].
 \label{eq:kl-contamination}
\end{equation}
If \(c=0\) under \(P\) and \(c=C_{\max}>0\) under \(R\),
compute the generated cost and both KL directions exactly.
\item For \(0<P(A)<1\), define the conditional law on \(\cX\) by
\[
 Q(x)=\frac{P(x)\I\{x\in A\}}{P(A)},\qquad x\in\cX.
\]
Prove
\begin{equation}
 \KL(Q,P)=-\log P(A),\qquad \KL(P,Q)=+\infty.
 \label{eq:kl-omission}
\end{equation}
What fraction of the target probability can be omitted with reverse KL
\(\log 2\)? If \(P\) is uniform on \(M\ge2\) zero-cost outcomes,
compute reverse KL and task cost for an arbitrary \(Q\) supported there.
\item Let \(B=\{x:c(x)>0\}\), \(P(B)=\varepsilon\in(0,1)\),
and \(Q(B)=q\). Prove \cref{eq:reverse-failure-prob}, and deduce an
upper bound on \(q\) when \(\KL(Q,P)\le K\).
For fixed \(q\in(0,1)\), compare the limits of
\(\operatorname{kl}(q,\varepsilon)\) and
\(\operatorname{kl}(\varepsilon,q)\) as \(\varepsilon\to0\).
Explain what changes when \(P(B)=0\).
\end{enumerate}
\begin{hints}
For (b), use the disjoint-support mixture calculation following
\cref{prop:positive-data}. For (c), apply data processing to
\(x\mapsto\I\{x\in B\}\), then use the entropy bound for a bit.
\end{hints}
\end{exercisechunk}

% Shared exercise chunk: l4-reverse-variational
\begin{exercisechunk}
\item \textbf{Variational representation of reverse KL.}
\label[exercise]{ex:reverse-variational}
Let \(P,Q\) be laws on a finite space \(\cX\), with \(X\sim P\)
and \(Y\sim Q\). Prove the identity in \cref{eq:reverse-adversarial}:
\[
 \KL(Q,P)=\sup_{g:\cX\to\R}
 \left\{1+\E[g(Y)]-\E[e^{g(X)}]\right\}.
\]
Include the cases where either law has zero probabilities.
When \(P,Q\) are strictly positive, show that the supremum is attained
by \(g^*(x)=\log(Q(x)/P(x))\).
\begin{hints}
Optimize each value \(g(x)\) separately.
\end{hints}
\end{exercisechunk}

% Shared exercise chunk: l4-stability
\begin{exercisechunk}
\item \textbf{Contraction and repeated errors.}
\label[exercise]{ex:stability}
This exercise compares recovery from an initial error with continual
correction of new errors.
\begin{enumerate}[(a),beginpenalty=10000]
\item Prove \cref{eq:dobrushin}.
\item For the recovery family \(Q_r\), compute the contraction factor
and use \cref{eq:stability-recursion} to recover the terminal and
count bounds in \cref{prop:recovery}. Verify that the recurrence
inequality is an equality here, giving the exact costs.
\item Now let \(Q_{\varepsilon,r}\), with \(X_0=0\), use
\[
 Q_{\varepsilon,r}(X_{t+1}=1\mid X_t=0)=\varepsilon,\qquad
 Q_{\varepsilon,r}(X_{t+1}=0\mid X_t=1)=r,
\]
where \(0<\varepsilon<1\), \(0\le r\le1\), and \(0\le t<T\).
For \(p_t=Q_{\varepsilon,r}(X_t=1)\), prove
\begin{equation}
 p_0=0,\qquad
 p_{t+1}=\varepsilon+(1-\varepsilon-r)p_t,\qquad
 p_t=\frac{\varepsilon}{\varepsilon+r}
       [1-(1-\varepsilon-r)^t].
 \label{eq:recovery}
\end{equation}
Identify the choices of \(r\) giving the two models in
\cref{ex:binary-count}.
\item Compute the contraction factor \(|1-\varepsilon-r|\).
For \(0\le r\le1-\varepsilon\), prove
\begin{equation}
 \E_{Q_{\varepsilon,r}}[c_{\mathrm{fin}}]\le
 \frac{\varepsilon}{\varepsilon+r},\qquad
 \frac1T\E_{Q_{\varepsilon,r}}[c_{\mathrm{cnt}}]\le
 \frac{\varepsilon}{\varepsilon+r}.
 \label{eq:recovery-cost}
\end{equation}
Recover these bounds from \cref{thm:markov-stability}, using
\(e_1=\varepsilon\) and local error \(\varepsilon\) at every
subsequent step. At \(r=0\), why does the positive contraction margin
still give the bound one? How large must \(r\) be relative to
\(\varepsilon\) for the limiting error probability to be at most
\(s\in(0,1)\)?
\end{enumerate}
\begin{hints}
For (a), normalize the positive and negative parts of
\(\alpha-\beta\) to obtain two probability laws.
\end{hints}
\end{exercisechunk}

% Shared exercise chunk: l4-cumulative-recovery
\begin{exercisechunk}
\item \textbf{Cumulative cost under recovery.}
\label[exercise]{ex:cumulative-recovery}
\label[exercise]{cor:cumulative-recovery}
Use the notation and assumptions of \cref{thm:markov-stability}, with
\(\rho<1\). For \(c_t:[N]\to[0,1]\), define
\[
 c_{\mathrm{sum}}(x_{1:T})=\sum_{t=1}^T c_t(x_t).
\]
Let \(X_{1:T}\sim P\) and \(Y_{1:T}\sim Q\). Prove
\begin{equation}
 \left|\E[c_{\mathrm{sum}}(X_{1:T})]
       -\E[c_{\mathrm{sum}}(Y_{1:T})]\right|
 \le\frac{e_1+\sum_{t=2}^T\varepsilon_t}{1-\rho}.
 \label{eq:cumulative-recovery}
\end{equation}
\begin{hints}
Sum the marginal bounds before bounding the geometric sums.
\end{hints}
\end{exercisechunk}

% Shared exercise chunk: l4-persistent-memory
\begin{exercisechunk}
\item \textbf{Forgetting errors and retaining information.}
\label[exercise]{ex:persistent-memory}
Use the reset mixture and persistent-state model from
\cref{sec:selective-memory}.
\begin{enumerate}[(a),beginpenalty=10000]
\item Prove the contraction bound for the reset mixture \(K_r\).
Construct a process whose next-token law uses only the last token
and whose initial bit can be recovered from every later token.
\item In \cref{eq:persistent-working-memory}, suppose the conditional
working-state kernels contract uniformly with factor \(\rho<1\).
Explain why contraction applies when comparing equal \(M\), while
different persistent \(M\)'s can remain distinguishable indefinitely.
\end{enumerate}

\end{exercisechunk}


\clearpage
\appendix
\renewcommand{\thesection}{4.\Alph{section}}

\section{Where finiteness is used}
\label{app:finiteness}

The finite sequence space guarantees bounded costs and permits
term-by-term differentiation. This section records which properties the
final set of results uses, so we can also apply them to other probability
spaces. The remaining appendices follow the order in which their topics first
require deferred arguments in the main text: reverse-KL concentration,
forward-KL variance control, finite-sample fitting, reverse-KL optimization,
and history-dependent recovery.

\paragraph{Distributional comparisons.}
Data processing, the TV bound (\cref{prop:tv-cost}), Pinsker's inequality,
the zero-cost bound (\cref{prop:positive-data}), the exponential-moment
bound (\cref{prop:entropy-cost}), and the variance-sensitive bounds
(\cref{prop:forward-variance,prop:hellinger-cost}) hold on general
probability spaces with the same bounded-cost assumptions. Their proofs
use integrals of densities in place of sums. The mixture and conditioning
identities in \cref{ex:kl-directions} also extend, with integrable costs.
The support requirement becomes \emph{absolute continuity}:
finite \(\KL(Q, P)\) requires
\(P(A)=0\Rightarrow Q(A)=0\) for every measurable event \(A\).
The Markov stability argument in \cref{thm:markov-stability} uses only
the stated contraction and expected local-error bounds. On general
state spaces, the same proof applies when the conditional laws exist.

\paragraph{Extrema at fixed KL.}
For fixed \(P\), each nonempty KL level set in
\cref{eq:fixed-kl-endpoints} is compact. Indeed, the probability simplex
is compact, and \(\KL(P, Q)\) is continuous wherever
\(Q(x)>0\) on the support of \(P\). It tends to infinity if one of
those probabilities tends to zero. Hence every finite level set is
closed. The infinite level set is the finite union of the closed faces
\(\{Q:Q(x)=0\}\) with \(P(x)>0\). Since \(d_c(P,Q)\) is
continuous in \(Q\), its minimum and maximum on each level are attained.
This justifies both endpoint formulas in \cref{prop:endpoints-moduli}.
On other spaces, compact level sets and continuity of the cost
discrepancy give the same conclusion. The budget definitions of the
moduli, their monotonicity, and the lower-bound inversion implication
use only well-defined discrepancies; identifying an exact converse
modulus from lower endpoints additionally uses attainment of those
endpoints.

\paragraph{Boundedness and exponential moments.}
Every real-valued cost on \([N]^T\) is bounded, so its exponential moments
exist for every \(\lambda\) and it has a finite subgaussian proxy.
Bounded costs on any probability space have these same properties.
For an unbounded cost, \cref{prop:entropy-cost} still holds when
\(\E_P[|c|],\E_Q[|c|]<\infty\), using those
\(\lambda>0\) for which \(\E_P[e^{\lambda c}]<\infty\).
The subgaussian condition in \cref{thm:reverse-cost} supplies the needed
exponential moments in both directions. Attainment of the optimal proxy
requires only that a finite valid proxy exist; \cref{app:optimal-proxy}
proves this directly.
The entropy decomposition in \cref{eq:population-kl} uses finite
\(H(P)\); for countable alphabets, \(H(P)<\infty\) keeps that
decomposition meaningful for comparing log loss with KL.

\paragraph{Differentiating expectations.}
Finite sums justify the reverse-KL gradient in \cref{app:reverse-optimization}.
For general densities,
sufficient conditions are differentiability on a fixed support and
integrable bounds, uniform in a neighborhood of the parameter, on the
derivatives of the density and the objective integrand. Dominated
convergence then permits differentiation under the integral.
For the exponential tilts in \cref{app:reverse-proof}, boundedness of
\(c\) supplies such bounds on every bounded interval of \(\lambda\).

\paragraph{The maximum-cost endpoint in the exact modulus.}
The proof of \cref{prop:reverse-modulus} uses finiteness to obtain a
maximum cost \(m\) with \(P(c=m)>0\). For a bounded cost on a general
space, take \(m\) to be its \emph{essential supremum} under \(P\):
the smallest number with \(P(c\le m)=1\).
If \(P(c=m)>0\), the same endpoint argument applies. If
\(P(c=m)=0\), the tilted KL increases to infinity: the tilted laws
concentrate on sets \(\{c>m-\varepsilon\}\) whose \(P\)-probabilities
decrease to zero. Every finite positive KL budget is therefore reached
by a finite tilt. Thus the exact-modulus identity still holds for
bounded costs; finiteness simplifies its endpoint proof.

\paragraph{Finite model classes.}
The finite-class guarantee in \cref{thm:finite-mle} uses a union bound
over models, producing \(\log|\mathcal Q|\); its likelihood-ratio
argument also works for general observation spaces with densities.
Finite-dimensional state representations are treated in Lecture~5.

\clearpage


\section{Reverse KL and exponential tilting}
\label{app:reverse-proof}

\subsection{From exponential moments to subgaussian control}

\entropycost*
\begin{proof}
The basic inequality uses the full exponential-moment function of the
cost. Put
\(\psi(\lambda)=\log\E_P[e^{\lambda(c-\mu)}]\).
Define the tilted law by
\[
 P_\lambda(x)=P(x)e^{\lambda(c(x)-\E_P\,c)-\psi(\lambda)}.
\]
This is a probability law. Writing \(\kappa=\KL(Q, P)\), we have
\[
 0\le\KL(Q, P_\lambda)
 =\kappa-\lambda(\E_Q[c]-\E_P[c])+\psi(\lambda).
\]
Consequently,
\[
 \E_Q[c]-\E_P[c]
 \le\inf_{\lambda>0}\frac{\kappa+\psi(\lambda)}{\lambda},
\]
which proves \cref{prop:entropy-cost}.
\end{proof}

Under the subgaussian condition,
\[
 \E_Q[c]-\E_P[c]
 \le\inf_{\lambda>0}
 \left(\frac{\kappa}{\lambda}+\frac{\lambda v^2}2\right)
 =v\sqrt{2\kappa}.
\]
Apply the same argument to \(-c\) for the other direction.
If \(v=0\), take \(\lambda\to\infty\) in each direction.

For completeness, the exponential-moment bound for averages follows
from Hoeffding's lemma. If \(Y\in[0,1]\), let
\(h(s)=\log\E[e^{s(Y-\E[Y])}]\).
Then \(h(0)=h'(0)=0\), and \(h''(s)\) is the variance of \(Y\)
under its tilted law. Every law on \([0,1]\) has variance at most
\(1/4\), giving \(h(s)\le s^2/8\).
For independent \(f_t(X_t)\in[0,1]\), factor the expectation of the
product (Lecture~2, \cref{lecture2-ex:probability-bounds}(a)) to obtain
\begin{align*}
 &\log\E\exp\left[
 \frac{\lambda}{T}\sum_{t=1}^T
 (f_t(X_t)-\E[f_t(X_t)])\right]\\
 &\qquad=\sum_{t=1}^T\log\E\exp\left[
 \frac{\lambda}{T}(f_t(X_t)-\E[f_t(X_t)])\right]\\
 &\qquad\le T\,\frac{(\lambda/T)^2}{8}
 =\frac{\lambda^2}{8T}.
\end{align*}
Hence the variance proxy is \(v^2=1/(4T)\).

\subsection{Optimal variance proxies}
\label{app:optimal-proxy}

The smallest valid proxy is attained whenever any finite proxy exists,
even on an infinite probability space. To see this, let \(c\) be an
integrable subgaussian cost and write
\(\psi(\lambda)=\log\E_P[e^{\lambda(c-\E_P[c])}]\).
The constraints defining a valid proxy give the identity
\[
 v_P(c)=\sup_{\lambda\ne0}\frac{2\psi(\lambda)}{\lambda^2}.
\]
This value satisfies every constraint, so the set of valid proxies is
exactly \([v_P(c),\infty)\).

For the bounded costs used here, expansion at zero gives
\[
 \psi(\lambda)=\frac{\lambda^2\sigma_P(c)^2}{2}+o(\lambda^2),
\]
and hence \(\sigma_P(c)^2\le v_P(c)\).
Scaling Hoeffding's lemma from \([0,1]\) to \([0,C_{\max}]\) gives
\(\psi(\lambda)\le\lambda^2C_{\max}^2/8\), so
\(v_P(c)\le C_{\max}^2/4\).

\subsection{The exact reverse-KL cost modulus}
\label{app:reverse-modulus}

How much can the expected cost increase if the reverse KL is at most
\(\kappa\)? Retaining the full exponential-moment function answers this
question exactly when we allow all model laws.

\begin{proposition}[Exact reverse-KL cost conversion]
\label{prop:reverse-modulus}
Fix a law \(P\) and a cost \(c\).
Let \(\psi(\lambda)=\log\E_P[
e^{\lambda(c-\E_P[c])}]\). For every \(\kappa\ge0\),
\begin{equation}
 \sup_{Q:\,\KL(Q, P)\le\kappa}
       (\E_Q[c]-\E_P[c])
 =\inf_{\lambda>0}\frac{\kappa+\psi(\lambda)}{\lambda}.
 \label{eq:reverse-modulus}
\end{equation}
The supremum ranges over all probability laws \(Q\).
\end{proposition}

\begin{proof}
\Cref{eq:entropy-cost} gives the upper bound. We construct laws attaining
it. Put \(\mu=\E_P[c]\) and use the tilted laws
\(P_\lambda(x)=P(x)e^{\lambda(c(x)-\mu)-\psi(\lambda)}\).
Differentiating \(\psi\) gives
\[
 \psi'(\lambda)=\E_{P_\lambda}[c]-\mu,\qquad
 \KL(P_\lambda, P)=\lambda\psi'(\lambda)-\psi(\lambda).
\]
If \(c\) is constant on \(\operatorname{supp}(P)\), both sides of
\cref{eq:reverse-modulus} are zero. Otherwise let
\[
 m=\max_{x:\,P(x)>0}c(x),\qquad
 \mathcal A=\{x:P(x)>0,\ c(x)=m\},\qquad
 \kappa_*=-\log P(\mathcal A)>0.
\]
As \(\lambda\) increases from zero to infinity, the tilted KL
increases continuously from zero to \(\kappa_*\). Indeed, its derivative
is \(\lambda\psi''(\lambda)>0\), and
\(P_\lambda\to P(\,\cdot\mid\mathcal A)\).
For \(0<\kappa<\kappa_*\), choose \(\lambda>0\) with
\(\KL(P_\lambda, P)=\kappa\). Then
\[
 \E_{P_\lambda}[c]-\mu
 =\psi'(\lambda)=\frac{\kappa+\psi(\lambda)}{\lambda},
\]
so this law attains the upper bound.

For \(\kappa\ge\kappa_*\), the law \(P(\,\cdot\mid\mathcal A)\) is feasible
and has cost \(m\), the largest possible cost on the support of \(P\).
Also \((\kappa+\psi(\lambda))/\lambda\to m-\mu\) as
\(\lambda\to\infty\), proving equality in this case.
For \(\kappa=0\), the only feasible law is \(P\), and
\(\psi(\lambda)/\lambda\to0\) as \(\lambda\downarrow0\).
\end{proof}

For the absolute discrepancy, apply the proposition to both \(c\)
and \(-c\):
\[
 \sup_{Q:\,\KL(Q, P)\le\kappa}d_c(P,Q)
 =\max_{s\in\{-1,1\}}\;
   \inf_{\lambda>0}\frac{\kappa+\psi(s\lambda)}{\lambda}.
\]
This is the reverse-KL counterpart of the forward-KL modulus in
\cref{eq:cost-modulus}. The subgaussian conversion replaces
\(\psi(s\lambda)\) by \(\lambda^2v_P(c)/2\). Any slack introduced
at that step comes from summarizing the cost's full exponential-moment
function by one quadratic bound.

\subsection{A full-support version of the binary comparison}

The infinite reverse KLs in the exact all-zero example are caused by
its support. A small perturbation gives finite values and still
distinguishes the generators.
Take \(P_\delta=\operatorname{Bernoulli}(\delta)^{\otimes T}\);
compare \(Q_{\rm ind}\) from \cref{ex:binary-count} with
\(Q_{\varepsilon,r}\) from \cref{ex:stability}. For \(0<\delta,\varepsilon,r<1\), all three
laws have full support. With
\(\operatorname{kl}(a,b)=a\log(a/b)+(1-a)\log((1-a)/(1-b))\),
the chain rule gives
\begin{align*}
 \KL(P_\delta, Q_{\varepsilon,r})
 &=\operatorname{kl}(\delta,\varepsilon)
 +(T-1)\bigl[(1-\delta)\operatorname{kl}(\delta,\varepsilon)
       +\delta\operatorname{kl}(\delta,1-r)\bigr],\\
 \KL(Q_{\varepsilon,r}, P_\delta)
 &=\sum_{t=1}^T
 \bigl[(1-p_{t-1})\operatorname{kl}(\varepsilon,\delta)
       +p_{t-1}\operatorname{kl}(1-r,\delta)\bigr].
\end{align*}
Here \(p_t\) is the generated probability of a one, given by
\cref{eq:recovery}. For the independent law the two KLs are
\(T\operatorname{kl}(\delta,\varepsilon)\) and
\(T\operatorname{kl}(\varepsilon,\delta)\).
At \(T=1000,\varepsilon=10^{-3},\delta=r=10^{-6}\):
\begin{center}
\begin{tabular}{lrrr}
\toprule
\(Q\) & \(\KL(P_\delta, Q)\)
 & \(\KL(Q, P_\delta)\) & \(\E_Q[c_{\mathrm{cnt}}]\)\\
\midrule
\(Q_{\rm ind}\) & \(0.993\) & \(5.909\) & \(1\)\\
\(Q_{\varepsilon,r}\) & \(1.006\) & \(5082.198\) & \(368.224\)\\
\bottomrule
\end{tabular}
\end{center}
The expected data cost is \(0.001\).
Forward KL weights the transition out of one by the tiny data
probability \(\delta\); reverse KL weights it by the much larger
generated probability \(p_{t-1}\).

\section{The variance-sensitive bound and comparisons}
\label{app:forward-cost}

\subsection{Cost variance and Hellinger distance}
\label{app:hellinger}

The \emph{Hellinger distance} between two laws \(P,Q\) is
\[
 h(P,Q)^2=\sum_x(\sqrt{P(x)}-\sqrt{Q(x)})^2.
\]
We use the normalization \(0\le h(P,Q)^2\le2\).
Expanding the square and using \(\sum_xP(x)=\sum_xQ(x)=1\) gives
\begin{equation}
 \sum_x\sqrt{P(x)Q(x)}=1-\frac{h(P,Q)^2}{2}.
 \label{eq:hellinger-affinity}
\end{equation}
The next bound uses this distance to control the cost discrepancy.

\begin{proposition}[Hellinger distance and cost variance]
\label{prop:hellinger-cost}
If \(0\le c\le C_{\max}\), then
\begin{equation}
 d_c(P,Q)\le2\sigma_P(c)h(P,Q)+C_{\max}h(P,Q)^2.
 \label{eq:hellinger-cost}
\end{equation}
Moreover, \(h(P,Q)^2\le\KL(P, Q)\).
\end{proposition}

\begin{proof}
Put \(\Delta(x)=\sqrt{Q(x)}-\sqrt{P(x)}\) and
\(\bar c=c-\E_P[c]\). Expand
\[
 \E_Q[c]-\E_P[c]
 =2\sum_x\bar c(x)\sqrt{P(x)}\Delta(x)
   +\sum_x\bar c(x)\Delta(x)^2.
\]
Cauchy--Schwarz bounds the first term in absolute value by
\(2\sigma_P(c)h(P,Q)\); \(|\bar c|\le C_{\max}\) bounds the second.

Let \(X\sim P\) and \(\kappa=\KL(P,Q)<\infty\).
Jensen's inequality for \(\exp\) gives
\begin{align*}
 \sum_x\sqrt{P(x)Q(x)}
 &=\E\left[\exp\left(-\frac12\log\frac{P(X)}{Q(X)}\right)\right]\\
 &\ge\exp\left(-\frac12\E\left[\log\frac{P(X)}{Q(X)}\right]\right)
 =e^{-\kappa/2}.
\end{align*}
Now \cref{eq:hellinger-affinity} and \(1-e^{-u}\le u\) give
\(h(P,Q)^2\le2(1-e^{-\kappa/2})\le\kappa\).
\end{proof}

Substituting \(h(P,Q)^2\le\KL(P, Q)\) into
\cref{eq:hellinger-cost} proves the following restated result.

\forwardvariance*
\begin{proof}
Apply \cref{eq:hellinger-cost} and \(h(P,Q)^2\le\KL(P,Q)\). Since
\(h(P,Q)\le\sqrt{\KL(P,Q)}\), the stated bound follows.
\end{proof}

\subsection{Which bound is stronger? Two single-bit comparisons}
\label{app:bound-comparison}

One bit suffices to show that either the variance-sensitive or the
subgaussian bound can be smaller. Throughout this subsection, \(c(x)=x\)
on \(\{0,1\}\), so \(C_{\max}=1\).

\begin{proposition}[Comparison near a fair coin and near zero cost]
\label{prop:bound-comparison}
\begin{enumerate}[(a)]
\item Suppose
\[
 P=\operatorname{Bernoulli}(1/2),\qquad
 Q=\operatorname{Bernoulli}(1/2+\delta),\qquad 0<|\delta|<1/2.
\]
As \(\delta\to0\), the Pinsker and
subgaussian bounds are asymptotic to \(d_c(P,Q)=|\delta|\), while
the variance-sensitive bound is asymptotic to \(\sqrt2|\delta|\).
\item Suppose
\[
 P=\operatorname{Bernoulli}(\varepsilon^2),\qquad
 Q=\operatorname{Bernoulli}(\varepsilon),\qquad 0<\varepsilon<1.
\]
As \(\varepsilon\downarrow0\), the actual discrepancy and the
variance-sensitive bound are both asymptotic to \(\varepsilon\).
The Pinsker bound is asymptotic to \(\sqrt{\varepsilon/2}\), and
the subgaussian bound is asymptotic to
\(\sqrt{\varepsilon/\log(1/\varepsilon)}\).
\end{enumerate}
\end{proposition}

\begin{proof}
For (a), direct expansion gives
\[
 \begin{aligned}
 \KL(P, Q)&=-\tfrac12\log(1-4\delta^2)
     =2\delta^2+O(\delta^4),\\
 \KL(Q, P)&=2\delta^2+O(\delta^4).
 \end{aligned}
\]
Also \(\sigma_P(c)=1/2\) and \(v_P(c)=1/4\): every variance proxy
is at least the variance, and Hoeffding's lemma gives equality for a
fair coin. Thus the reverse-KL term of
\cref{eq:two-concentration-directions} is asymptotic to
\(|\delta|\). The minimum of its two terms is between the actual
discrepancy and this term, so it has the same asymptotic value.
Substitution into \cref{eq:pinsker-cost,eq:forward-variance} gives
the other two comparisons.

For (b), the Bernoulli KL formula gives
\[
 \begin{aligned}
 d_c(P,Q)&=\varepsilon-\varepsilon^2\sim\varepsilon,\\
 \KL(P, Q)&\sim\varepsilon,\qquad
 \KL(Q, P)\sim\varepsilon\log(1/\varepsilon),\\
 \sigma_P(c)&=\varepsilon\sqrt{1-\varepsilon^2}\sim\varepsilon.
 \end{aligned}
\]
The variance-sensitive bound is therefore
\(2\sigma_P(c)\sqrt{\KL(P, Q)}+\KL(P, Q)
\sim\varepsilon\).
To evaluate the subgaussian bound, we use the sharp Bernoulli proxy
\citep[Eq.~(1.1) and Theorem~1.1]{schlemm2014}:
\[
 v_{\operatorname{Bernoulli}(p)}(c)
 =\frac{1-2p}{2\log((1-p)/p)},\qquad
 0<p<1,\quad p\ne\tfrac12,
\]
with value \(1/4\) at \(p=1/2\).
In particular,
\[
 v_Q(c)\sim\frac1{2\log(1/\varepsilon)},\qquad
 v_P(c)\sim\frac1{4\log(1/\varepsilon)}.
\]
Consequently,
\[
 \begin{aligned}
 \sqrt{2v_Q(c)\KL(P, Q)}
    &\sim\sqrt{\frac{\varepsilon}{\log(1/\varepsilon)}},\\
 \sqrt{2v_P(c)\KL(Q, P)}
    &\sim\sqrt{\frac{\varepsilon}{2}}.
 \end{aligned}
\]
The forward-KL term is the smaller one for sufficiently small
\(\varepsilon\). Pinsker gives \(\sqrt{\KL(P, Q)/2}
\sim\sqrt{\varepsilon/2}\).
\end{proof}

Near a fair coin, Pinsker and the subgaussian inequality give the
correct leading constant. Near zero cost, the variance-sensitive bound
uses the small data variance to give the correct leading order, whereas
even the optimal subgaussian proxies give a much larger bound.
The first example requires square-root dependence on KL in a general
cost bound. The zero-cost equality example in
\cref{sec:kl-cost} requires linear dependence when the data variance
vanishes. These explain the two terms in \cref{eq:forward-variance}.

\subsection{Moduli of continuity in analysis and approximation}
\label{app:moduli-analysis}

A modulus of continuity quantifies how much a function's output can
change when its input changes by at most a specified amount. For a
function \(h:[0,1]\to\R\), its \emph{modulus of continuity} is
\[
 \omega_h(\delta)
 =\sup_{\substack{x,y\in[0,1]\\|x-y|\le\delta}}|h(x)-h(y)|.
\]
Instead of only asking whether nearby inputs have nearby outputs, we
ask how nearby the outputs must be. For example, a Lipschitz bound
\(|h(x)-h(y)|\le L|x-y|\) gives \(\omega_h(\delta)\le L\delta\).

\begin{proposition}[Quantitative uniform continuity]
A function \(h:[0,1]\to\R\) is uniformly continuous if and only if
\(\omega_h(\delta)\to0\) as \(\delta\downarrow0\).
\end{proposition}

\begin{proof}
Uniform continuity makes every difference in the defining supremum at
most \(\varepsilon\) once \(\delta\) is sufficiently small. Conversely,
\(\omega_h(\delta)<\varepsilon\) gives
\(|h(x)-h(y)|<\varepsilon\) whenever \(|x-y|\le\delta\).
\end{proof}

In our comparison, the function is the expected-cost map
\(P\mapsto\E_P[c]\). We measure input discrepancy using KL and
output discrepancy using \(d_c\). Taking the supremum over \(P\)
as well as \(Q\) gives the uniform comparison. If \(\mathcal P\)
contains all data laws, Pinsker's inequality gives
\(\omega^{\mathcal P}_{D\to c}(\kappa)\le C_{\max}\sqrt{\kappa/2}\) for
\(0\le c\le C_{\max}\): small KL uniformly forces nearby expected costs.

The expected-cost map can send distinct laws to the same value.
Accordingly, \(d_c\) is a \emph{pseudometric}: it is symmetric,
satisfies the triangle inequality, and vanishes on equal laws, but can
also vanish on distinct laws. The converse modulus quantifies how
far apart the input laws can be while their expected costs remain
close. \Cref{prop:binary-moduli,prop:terminal-moduli} show how much
this depends on the cost and reference law.

\paragraph{From continuity to approximation rates.}
A quantitative continuity bound can also tell us how accurately a
function can be approximated by polynomials. Jackson's theorem states
that, for a universal constant \(C\), every continuous
\(h:[0,1]\to\R\) satisfies
\[
 \inf_{\deg p\le n}\|h-p\|_\infty
 \le C\,\omega_h(1/n),\qquad n\ge1.
\]
Thus \(\omega_h(\delta)\le L\delta^\alpha\), with
\(0<\alpha\le1\), guarantees a degree-\(n\) approximation with
uniform error at most \(CLn^{-\alpha}\). The modulus turns a
qualitative continuity statement into a bound on the degree needed
for a given approximation error.

\section{Likelihood fitting and conditional support}
\label{app:learning}

\subsection{Proof of the finite-class likelihood guarantee}
\label{app:finite-mle}

\finitemle*
\begin{proof}
We first prove \(h(P,\widehat Q)^2\le b\) with probability at least
\(1-\delta\), then apply \cref{prop:hellinger-cost}.
The probability tools below are proved in Lecture~2,
\cref{lecture2-ex:probability-bounds}; the inequality
\(1-u\le e^{-u}\) is proved in
\cref{lecture2-ex:sample-size-bounds}(a) of that lecture.

For a fixed \(Q\in\mathcal Q\), define the likelihood ratio
\[
 \Lambda_Q=\prod_{i=1}^n
 \frac{Q(X^{(i)}_{1:T})}{P(X^{(i)}_{1:T})}.
\]
The denominators are positive almost surely because each observed
sequence has law \(P\). Independence of the sequences gives
\begin{align*}
 \E[\sqrt{\Lambda_Q}]
 &=\prod_{i=1}^n\E\left[
       \sqrt{\frac{Q(X^{(i)}_{1:T})}{P(X^{(i)}_{1:T})}}\right]\\
 &=\left(\sum_{x:P(x)>0}P(x)\sqrt{\frac{Q(x)}{P(x)}}\right)^n
 =\left(\sum_x\sqrt{P(x)Q(x)}\right)^n.
\end{align*}
Using \cref{eq:hellinger-affinity} and \(1-u\le e^{-u}\),
\[
 \E[\sqrt{\Lambda_Q}]
 =\left(1-\frac{h(P,Q)^2}{2}\right)^n
 \le e^{-nh(P,Q)^2/2}.
\]

Since \(P\in\mathcal Q\), \cref{eq:approximate-fit} gives
\(\widehat L(\widehat Q)\le\widehat L(P)+\eta\). Hence
\[
 \log\Lambda_{\widehat Q}
 =nT\bigl(\widehat L(P)-\widehat L(\widehat Q)\bigr)
 \ge-nT\eta.
\]
For a fixed candidate \(Q\), apply Markov's inequality to
\(\sqrt{\Lambda_Q}\ge0\):
\begin{align*}
 \P(\Lambda_Q\ge e^{-nT\eta})
 &=\P(\sqrt{\Lambda_Q}\ge e^{-nT\eta/2})\\
 &\le e^{nT\eta/2}\E[\sqrt{\Lambda_Q}]\\
 &\le\exp\left(\frac{nT\eta}{2}-\frac{nh(P,Q)^2}{2}\right).
\end{align*}
For every \(Q\) in the set
\(\mathcal B=\{Q\in\mathcal Q:h(P,Q)^2>b\}\), this probability is
at most \(\delta/|\mathcal Q|\). The selected model can belong to
\(\mathcal B\) only if one of these likelihood-ratio events occurs:
\[
 \{h(P,\widehat Q)^2>b\}
 \subseteq\bigcup_{Q\in\mathcal B}
          \{\Lambda_Q\ge e^{-nT\eta}\}.
\]
The union bound therefore yields
\[
 \P(h(P,\widehat Q)^2>b)
 \le\sum_{Q\in\mathcal B}\P(\Lambda_Q\ge e^{-nT\eta})
 \le |\mathcal B|\frac{\delta}{|\mathcal Q|}\le\delta.
\]
On the complementary event, \cref{prop:hellinger-cost} gives
\(d_c(P,\widehat Q)\le2\sigma_P(c)\sqrt b+C_{\max}b\).
If \(\E_P[c]=0\), nonnegativity of \(c\) implies that \(c=0\)
\(P\)-almost surely, so \(\sigma_P(c)=0\) and
\(d_c(P,\widehat Q)=\E_{\widehat Q}[c]\).
\end{proof}

The independent observations here are complete sequences. Dependence
between tokens is part of their joint law \(P\); it never enters this
likelihood-ratio calculation as an independence assumption.
Lecture~6 studies how much information dependent observations supply.

\subsection{Proofs of the parameter-sharing results}
\label{app:growing-horizon}

\paragraph{Two-sided risk and the leading constant.}
\sharingrisk*

\begin{proof}[Proof of \cref{prop:sharing-kl}]
First consider \(B\sim\operatorname{Binomial}(L,p)\), with
\(L\ge1\) and \(0<p<1\).
The smoothed estimator is \(\widehat p=(B+1)/(L+2)\).
Using \(1/(B+1)=\int_0^1u^B\,du\),
\[
 \E\!\left[\frac1{B+1}\right]
 =\int_0^1(1-p+pu)^L\,du
 =\frac{1-(1-p)^{L+1}}{(L+1)p}.
\]
Consequently,
\[
 p\,\E[1/\widehat p]\le\frac{L+2}{L+1},
 \qquad
 (1-p)\,\E[1/(1-\widehat p)]\le\frac{L+2}{L+1}.
\]
The second inequality follows by applying the same calculation to
\(L-B\). Jensen's inequality gives
\begin{align*}
 \E[\kl(p,\widehat p)]
 &\le p\log\!\bigl(p\,\E[1/\widehat p]\bigr)
 +(1-p)\log\!\bigl((1-p)\,\E[1/(1-\widehat p)]\bigr)\\
 &\le\log\frac{L+2}{L+1}\le\frac1{L+1}.
\end{align*}
For \(p=0,1\), the loss equals \(\log((L+2)/(L+1))\)
deterministically, so the bound also holds.

For group \(j\), the pooled count is binomial with \(L=nm_j\).
Independence of the test tokens makes the sequence KL the sum of
the Bernoulli KLs. Therefore
\[
 \kappa_n(T)
 =\sum_{j=1}^{d(T)}m_j\,\E[\kl(p_j,\widehat p_j)]
 \le\sum_{j=1}^{d(T)}\frac{m_j}{nm_j+1}
 \le\frac{d(T)}n.
\]
For the lower bound, the binomial mean and variance give
\[
 \E[(\widehat p-p)^2]
 =\frac{Lp(1-p)+(1-2p)^2}{(L+2)^2}.
\]
Pinsker's inequality implies
\[
 \E[\kl(p,\widehat p)]
 \ge2\E[(\widehat p-p)^2]
 \ge\frac{2Lp(1-p)}{(L+2)^2}.
\]
With \(L=nm_j\) and \(p_j\in[a,1-a]\),
\[
 m_j\,\E[\kl(p_j,\widehat p_j)]
 \ge\frac{2a(1-a)}n\left(\frac{nm_j}{nm_j+2}\right)^2
 \ge\frac{2a(1-a)}{9n}.
\]
Summing over groups proves \cref{eq:sharing-rate}, with explicit bounds
\[
 \frac{2a(1-a)}9\frac{d(T)}n
 \le\kappa_n(T)\le\frac{d(T)}n.
\]

For the leading constant, the same mean-squared error formula gives,
uniformly for \(p\in[a,1-a]\),
\[
 \E[(\widehat p-p)^2]=\frac{p(1-p)}L+O(L^{-2}).
\]
In a fixed neighborhood of \(p\), Taylor's theorem gives
\[
 \kl(p,q)=\frac{(q-p)^2}{2p(1-p)}+O(|q-p|^3).
\]
Binomial moment bounds give
\(\E[|\widehat p-p|^3]=O(L^{-3/2})\).
Outside that neighborhood the probability is exponentially small
by a binomial tail bound, while
\(\kl(p,\widehat p)\le\log(L+2)\).
Splitting the expectation between the two events yields
\[
 \E[\kl(p,\widehat p)]
 =\frac1{2L}+O(L^{-3/2}).
\]
The constants can be chosen uniformly for \(p\in[a,1-a]\),
\(0<a<1/2\).

Summing over groups gives
\[
 \kappa_n(T)=\frac{d(T)}{2n}
 +O_a\!\left(\frac1{n^{3/2}}
              \sum_{j=1}^{d(T)}m_j^{-1/2}\right).
\]
The relative error is \(O_a((n\min_j m_j)^{-1/2})\), which tends
to zero under the stated condition.
This proves \cref{eq:sharing-leading}.
\end{proof}

\paragraph{Learning rates and task-cost intervals.}
\sharingfullcost*

\begin{proof}
Let the \(U_t\) be independent fair bits and \(V_t=0\), giving one law
\(P\) on infinite sequences. Use terminal cost \(V_T\) and counting
cost \(\sum_{t=1}^T V_t\).
Since \(1\le\lceil d(T)\rceil\le T\), at each horizon we can partition
the first \(T\) positions into \(\lceil d(T)\rceil\) nonempty groups.
The estimator may use a different partition at each horizon; the data
law \(P\) remains fixed.
Let \(\widehat P^U_T\) be the smoothed pooled estimator and set
\(r_T=\E[\KL(P^U_T,\widehat P^U_T)]\). For every law \(R\)
supported on \(v_{1:T}\ne0^T\), define
\[
 \widehat Q_R=e^{-r_T}(\widehat P^U_T\otimes\delta_{0^T})
 +(1-e^{-r_T})R.
\]
Since all Bernoulli parameters are
fixed at \(1/2\), \(r_T\) is a deterministic function of the known
sample sizes. Thus the prescribed mixture weight defines an estimator
from the observed sample.

For this construction, \(P_T=P^U_T\otimes\delta_{0^T}\).
Since \(R\) assigns no mass to \(v_{1:T}=0^T\),
\[
 \KL(P_T,\widehat Q_R)
 =\KL(P^U_T,\widehat P^U_T)+r_T.
\]
Hence every \(\widehat Q_R\) has expected sequence KL
\(\kappa_n(T)=2r_T\asymp \lceil d(T)\rceil/n\asymp d(T)/n\).
Both costs vanish under the first
mixture component, so
\[
 \E\!\left[\E_{\widehat Q_R}[c]\right]
 =(1-e^{-r_T})\E_R[c].
\]
On \(v_{1:T}\ne0^T\), terminal cost ranges from zero to one, while count
cost ranges from one to \(T\). Varying \(R\) over point masses and their
mixtures gives the two intervals in \cref{prop:sharing-full-cost}, since
\(r_T=\kappa_n(T)/2\).
If an infinite-sequence output law is required, extend each fitted
length-\(T\) law by a fixed continuation.
\end{proof}

\subsection{A sharp statistical obstruction to reverse-KL fitting}
\label{app:conditional-support}

Suppose the learner must answer prompts drawn from a fixed law \(\mu\).
Replaying a few training prompts would change the task: the distribution
of prompts is supplied externally. This makes support recovery a
conditional prediction problem.

\begin{proposition}[Two candidate rules and rare distinguishing prompts]
\label{prop:conditional-support}
Let \(\mathcal U\) be finite, with \(\mu(u)>0\) for every \(u\).
Two known functions \(f_0,f_1:\mathcal U\to\{0,1\}\) disagree on
\(\mathcal A\subseteq\mathcal U\), where \(\mu(\mathcal A)=\alpha>0\).
Under candidate \(i\in\{0,1\}\), a demonstration has law
\[
 P_i(u,y)=\mu(u)\I\{y=f_i(u)\}.
\]
After \(n\) independent demonstrations, a possibly randomized learner
returns a conditional law \(\widehat q\), inducing
\(\widehat Q(u,y)=\mu(u)\widehat q(y\mid u)\).
Then
\begin{equation}
 \inf_{\mathcal L}\max_{i\in\{0,1\}}
 \mathbb P_{P_i^{\otimes n},\mathcal L}
 \bigl(\KL(\widehat Q, P_i)=+\infty\bigr)
 =\frac12(1-\alpha)^n.
 \label{eq:support-lower}
\end{equation}
Here \(\mathcal L\) ranges over all such learning rules.
\end{proposition}

\begin{proof}
Finite reverse KL requires
\(\widehat q(f_i(u)\mid u)=1\) for every \(u\); hence it requires
\(\widehat Q=P_i\).
With probability \((1-\alpha)^n\), none of the training prompts lies
in \(\mathcal A\). Conditional on this event the data have the same
distribution under both candidates.
No output can equal both \(P_0\) and \(P_1\).
Averaging over the two equally likely candidates, the probability of
failure on this event is therefore at least \(1/2\).
This gives the lower bound on the worst-case failure probability.

For equality, use any observed prompt in \(\mathcal A\) to identify \(i\)
and return \(P_i\). If none occurs, choose one of the two candidates
with an independent fair coin. This has the stated failure probability
under each candidate.
\end{proof}

For \(0<\alpha<1\) and \(0<\delta<1/2\), success probability at least
\(1-\delta\) requires
\[
 n\ge\frac{\log(1/(2\delta))}{-\log(1-\alpha)}.
\]
To make this concrete, a prompt is a pair of \(m\)-bit integers
\((L,R)\), drawn uniformly and independently. The candidate rules are
\[
 f_0(L,R)=\I\{L\ge R\},\qquad f_1(L,R)=\I\{L>R\}.
\]
They differ exactly when \(L=R\), an event of probability \(2^{-m}\).
The sample requirement is consequently
\(\Omega(2^m\log(1/\delta))\) for \(0<\delta\le1/4\).
The class has only two rules; random demonstrations rarely distinguish
them.

For comparison, let \(q\) give the common answer off \(\mathcal A\)
and a fair bit on \(\mathcal A\). With \(Q(u,y)=\mu(u)q(y\mid u)\)
and task cost \(c_i(u,y)=\I\{y\ne f_i(u)\}\), direct calculation gives
\[
 \KL(P_i, Q)=\alpha\log2,\qquad
 \E_{P_i}[c_i]=0,\qquad
 \E_Q[c_i]=\frac\alpha2,\qquad
 \KL(Q, P_i)=+\infty.
\]
Finite reverse KL demands exact conditional support even on very rare
prompts. Forward KL and expected error can already be small before those
prompts have been observed.
One expert answer to a \emph{chosen} prompt with \(L=R\) identifies the
rule. This specifies additional information that resolves the difficulty.
Lectures~19 and~24 return to such differences between random examples
and targeted queries.

\clearpage
\section{Reverse-KL optimization}
\label{app:reverse-optimization}

Direct reverse-KL optimization uses target scores on generated samples.
We derive its gradient when these scores are supplied, then connect a
cost-defined target to KL-regularized reinforcement learning.

\subsection{A target known up to normalization}

Suppose a target law is specified as
\[
 P(x)=\frac{w(x)}{\mathcal Z},\qquad
 \mathcal Z=\sum_x w(x)>0,
\]
where \(w\) is evaluable but \(\mathcal Z\) is expensive to compute.
For \(\widetilde X\sim Q_\theta\), with \(Q_\theta\) supported where
\(w>0\),
\begin{equation}
 \KL(Q_\theta, P)
 =\E[\log Q_\theta(\widetilde X)-\log w(\widetilde X)]
   +\log\mathcal Z.
 \label{eq:unnormalized-target}
\end{equation}
The final term is independent of \(\theta\). We can sample from
\(Q_\theta\), evaluate \(\log Q_\theta-\log w\) on those samples,
and estimate gradients without sampling from \(P\).

For a differentiable model with fixed support contained in that of
\(w\), the gradient is
\[
 \E\!\left[
 \bigl(\log Q_\theta(\widetilde X)-\log w(\widetilde X)\bigr)
 \nabla_\theta\log Q_\theta(\widetilde X)
 \right],\qquad \widetilde X\sim Q_\theta.
\]
Here \(\widetilde X\) denotes the whole generated sequence.
Differentiating the expectation gives this identity; the additional
normalization derivative vanishes because \(\sum_x\nabla Q_\theta(x)=0\).
Monte Carlo averages of these gradients give stochastic optimization
updates. Eliminating the normalizer preserves the objective's minimizers;
it does not certify that optimization has reached a small KL.

\subsection{A cost-defined target connects reverse KL to RL}

Suppose instead that a cost \(c(x)\) can be evaluated.
Given a positive reference law \(R\) and \(\beta>0\), define
\[
 P_\beta(x)=\frac{R(x)e^{-\beta c(x)}}{\mathcal Z_\beta},
 \qquad \mathcal Z_\beta=\E[e^{-\beta c(X)}],\quad X\sim R.
\]
For \(Y\sim Q\), direct substitution gives
\begin{equation}
 \KL(Q, P_\beta)
 =\KL(Q, R)+\beta\E[c(Y)]+\log\mathcal Z_\beta.
 \label{eq:gibbs-cost}
\end{equation}
Thus minimizing reverse KL to \(P_\beta\) is equivalent to minimizing
\(\E[c(Y)]+\beta^{-1}\KL(Q,R)\).
The reference penalizes departures from an existing model, while the
cost evaluates generated behavior. The supplied cost defines a new
target whose relative probabilities can be evaluated. This is a useful
connection to KL-regularized reinforcement learning and post-training.

\section{Recovery for history-dependent generators}
\label{app:history-recovery}

Final-answer accuracy can be controlled by bounding the influence of an
earlier token on the final marginal. Markov contraction gives geometric
decay of this influence. Here we allow the generator to use the entire
prefix and quantify the influence directly.

\subsection{From the influence of a token to task cost}

Let \(Q\) be a generator with next-token laws specified at every prefix.
For a prefix \(h=x_{<t}\) followed by a supplied token \(x\), write
\(Q(X_s\in\cdot\mid h,x)\) for the law at position \(s\ge t\)
obtained by continuing this generator from \((h,x)\).
At \(s=t\) this law is \(\delta_x\).
Suppose that changing the supplied token changes later marginals by at most
\begin{equation}
 \sup_{h,x,x'}
 \tv{Q(X_s\in\cdot\mid h,x)-Q(X_s\in\cdot\mid h,x')}
 \le b_{s-t},\qquad s\ge t,\qquad b_0=1.
 \label{eq:history-influence}
\end{equation}
The coefficient \(b_j\ge0\) bounds the influence \(j\) positions later.
For a Markov generator satisfying \cref{eq:markov-contraction},
we can take \(b_j=\rho^j\).

For an arbitrary data law \(P\), let
\begin{equation}
 \varepsilon_t=
 \E_P\!\left[
 \tv{P(\cdot\mid X_{<t})-Q(\cdot\mid X_{<t})}
 \right],\qquad 1\le t\le T,
 \label{eq:history-local-error}
\end{equation}
where the prefix at \(t=1\) is empty.

\begin{proposition}[Recovery with history-dependent conditionals]
\label{prop:history-recovery}
Under \cref{eq:history-influence}, for \(c,c_t:[N]\to[0,1]\),
\begin{align}
 \left|\E_P[c(X_T)]-\E_Q[c(X_T)]\right|
 &\le\sum_{t=1}^T b_{T-t}\varepsilon_t,
 \label{eq:history-terminal}\\
 \left|\E_P\!\left[\sum_{s=1}^T c_s(X_s)\right]
       -\E_Q\!\left[\sum_{s=1}^T c_s(X_s)\right]\right|
 &\le\sum_{t=1}^T\varepsilon_t\sum_{j=0}^{T-t}b_j.
 \label{eq:history-cumulative}
\end{align}
In particular, if \(\sum_{j\ge0}b_j\le B\), the cumulative-cost
discrepancy is at most \(B\sum_{t=1}^T\varepsilon_t\).
\end{proposition}

\begin{proof}
Replace the generator's conditionals by the data conditionals one at
a time. Let \(R^{(t)}\) use \(P\) for the first \(t\) tokens and \(Q\)
thereafter:
\[
 R^{(t)}(x_{1:T})
 =P(x_{\le t})\prod_{j=t+1}^T Q(x_j\mid x_{<j}),
 \qquad R^{(0)}=Q,\quad R^{(T)}=P.
\]
The laws \(R^{(t)}\) and \(R^{(t-1)}\) have the same prefix distribution
\(P(X_{<t}\in\cdot)\). At a common prefix \(h\), their next-token laws
are \(P(\cdot\mid h)\) and \(Q(\cdot\mid h)\).
A \emph{coupling} puts two random variables with prescribed laws on
one probability space. On a finite set we can couple these tokens
to disagree with probability exactly their TV distance: assign
probability \(\min\{P(x\mid h),Q(x\mid h)\}\) to matching token \(x\),
then couple the remaining masses, whose supports are disjoint.

If the tokens agree, their \(Q\)-continuations have the same law.
If they differ, \cref{eq:history-influence} and the TV cost bound
limit the difference in expected cost at position \(s\ge t\)
to \(b_{s-t}\). Averaging over the coupled tokens and the data prefix gives
\[
 \left|\E_{R^{(t)}}[c_s(X_s)]
       -\E_{R^{(t-1)}}[c_s(X_s)]\right|
 \le b_{s-t}\varepsilon_t.
\]
For \(t>s\), this difference is zero.
Telescoping over \(t\) with \(s=T\) proves
\cref{eq:history-terminal}. Summing the same bounds over \(s\) and
reversing the two finite sums proves \cref{eq:history-cumulative}.
\end{proof}

An initial discrepancy contributes \(b_{T-1}\varepsilon_1\) to
the final-answer bound. Thus \(b_j\to0\) expresses recovery from
an initial error. Summability gives a bounded total contribution
to cumulative cost. With fresh errors \(\varepsilon_t\le\varepsilon\),
the final-answer bound is at most \(B\varepsilon\).

\subsection{A sufficient condition on next-token laws}

We can obtain the influence bounds by limiting sensitivity to each
past position. Let \(a_j\ge0\) bound the effect of a token \(j\) positions
back, and suppose that for every \(t\) and pair of prefixes,
\begin{equation}
 \tv{Q(\cdot\mid x_{<t})-Q(\cdot\mid y_{<t})}
 \le\sum_{j=1}^{t-1}a_j\mathbf 1\{x_{t-j}\ne y_{t-j}\},
 \qquad \sum_{j\ge1}a_j=\rho<1.
 \label{eq:history-sensitivity}
\end{equation}
There can be nonzero coefficients at arbitrarily large lags.

\begin{proposition}[Summable sensitivity gives bounded error propagation]
\label{prop:summable-sensitivity}
Under \cref{eq:history-sensitivity}, \cref{eq:history-influence} holds with
\begin{equation}
 b_0=1,\qquad
 b_k=\sum_{j=1}^k a_jb_{k-j},\qquad
 \sum_{k\ge0}b_k=\frac1{1-\rho}.
 \label{eq:influence-renewal}
\end{equation}
\end{proposition}
\begin{proof}
Start two \(Q\)-continuations from the same prefix with different
tokens at position \(t\). Couple each successive pair of next-token
laws to minimize disagreement, as in the previous proof.
Let \(d_k\) be the probability that the tokens at position \(t+k\)
differ. The supplied tokens give \(d_0=1\).
The earlier prefixes agree, so \cref{eq:history-sensitivity} gives
\[
 d_k\le\sum_{j=1}^k a_jd_{k-j}.
\]
Induction gives \(d_k\le b_k\). The TV distance of the two marginals
is at most their disagreement probability, proving
\cref{eq:history-influence}.
For \(B_m=\sum_{k=0}^m b_k\), the recurrence gives
\(B_m\le1+\rho B_m\), hence \(B_m\le1/(1-\rho)\).
The full sum is therefore finite. Summing the recurrence and
reordering nonnegative terms gives
\[
 \sum_{k\ge0}b_k=1+\rho\sum_{k\ge0}b_k=\frac1{1-\rho}.
\]
\end{proof}

The recurrence includes influence transmitted through intervening tokens:
\(b_1=a_1\), \(b_2=a_2+a_1^2\), and so on.
For first-order Markov contraction, \(a_1=\rho\) and \(a_j=0\)
for \(j>1\), giving \(b_k=\rho^k\).

\paragraph{Return to the KL question.}
For \(P=\delta_{0^T}\), the local errors in
\cref{eq:history-local-error} satisfy
\[
 \varepsilon_t=1-Q(0\mid0^{t-1})
 \le-\log Q(0\mid0^{t-1}).
\]
The KL chain rule, \cref{eq:population-kl}, therefore gives
\(\sum_t\varepsilon_t\le\KL(P,Q)\).
Under \cref{eq:history-sensitivity}, we obtain
\begin{equation}
 \E_Q[c_{\mathrm{fin}}]
 \le\sum_{t=1}^T b_{T-t}[-\log Q(0\mid0^{t-1})],
 \qquad
 \E_Q[c_{\mathrm{cnt}}]\le\frac{\KL(P,Q)}{1-\rho}.
 \label{eq:stable-zero-cost}
\end{equation}
For the final answer, recovery weights the local log losses by their
remaining influence. For the error count, a fixed contraction margin
removes the horizon factor in the unrestricted example
\(\E_{Q_+}[c_{\mathrm{cnt}}]=T(1-e^{-\kappa})\).

\subsection{Weak dependence at each step can preserve an old error}

A small effect on each prediction can still accumulate if the same
old error keeps affecting future predictions.
\begin{example}[Repeated use of the first token]
\label{ex:unchanging-influence}
Let \(P=\delta_{0^T}\), \(T\ge2\), \(0<\rho<1\), and
\[
 Q(X_1=1)=1-e^{-\kappa},\qquad
 Q(X_t=1\mid X_{<t})=\rho X_1,\quad t\ge2.
\]
At every time \(t\ge2\), any two next-token laws differ in TV by at
most \(\rho\). Nevertheless,
\[
 \KL(P,Q)=\kappa,\qquad
 \E_Q[c_{\mathrm{fin}}]=\rho(1-e^{-\kappa}),\qquad
 \E_Q[c_{\mathrm{cnt}}]=(1-e^{-\kappa})[1+\rho(T-1)].
\]
Indeed, the all-zero path has probability \(e^{-\kappa}\), and
\(Q(X_t=1)=\rho Q(X_1=1)\) at every \(t\ge2\).
Changing the first token keeps changing all later marginals by
\(\rho\), so \cref{eq:history-influence} requires \(b_j\ge\rho\)
for every \(j\ge1\) when the model is continued indefinitely.
The influences are not summable: the initial error keeps causing
mistakes instead of being forgotten.
\end{example}

\end{document}
